\chapter{数学基础与后续研究}
\label{chap:mathematics-and-future-research}

设一项研究已经把训练损失降到零，并发现参数更新趋于稳定。它是否已经证明模型在
新数据上可靠？若模型随后用于控制加热器，预测准确又是否足以保证温度不越界？
前一个问题需要比较数据选择后的总体风险，后一个问题需要检查行动作用下的整条轨迹。
计算成功、统计证据和行为保证针对不同对象，不能用一个结论包办。

前述各章已经建立函数、空间、随机变量、图、信号以及状态和策略的数学语言。
本章把这些对象放回具体研究任务：数据选出了什么候选，某个算子保留了什么关系，
计算返回的是导数、最优解、条件分布还是随机样本，证据又在哪个总体或决策制度下有效。
同一组数据和参数可以参与几项任务；每换一个问题，都需要重新说明比较范围和成立条件。

投影、谱和固定点尤其容易掩盖这种变化。最小二乘投影比较有限响应向量，
条件期望投影比较在某种信息下可用的随机变量；协方差谱比较方差方向，
状态转移谱考察反复作用；折扣价值的固定点与相互最佳响应的固定点也有不同的
空间和映射条件。表~\ref{tab:future-reused-mathematical-objects}把这些差别并列，
正文将用可复算的拟合、查询、决策和迁移例子逐项检验。

\begin{table*}[htbp]
  \centering
  \begin{tabularx}{\linewidth}{lXX}
    \toprule
    共用工具 & 当前对象与条件 & 应当读取的结论\\
    \midrule
    正交投影 & 最小二乘把有限响应向量投到设计矩阵列空间；条件期望把平方可积随机变量投到给定信息的可测子空间 & 前者残差与每列设计正交；后者残差与每个平方可积、已知信息可测的随机变量正交。投影空间不同\\
    特征值 & 协方差为对称半正定矩阵；状态转移矩阵可以非正规 & 协方差谱描述方差方向；转移谱描述长期模式，单靠特征值不能排除有限时间放大\\
    固定点 & 折扣Bellman算子在上确界范数下压缩；Nash最佳响应通常是集合值映射 & 前者给唯一价值与迭代收敛；后者在相应条件下给均衡存在，不能直接推出唯一性或学习动态收敛\\
    沿轨迹的标量变化 & Lyapunov函数正定且不增；最优价值满足成本递推；屏障函数规定允许集合及可行动作 & 分别检验稳定、最优成本和逐轨迹安全。成本低或最终收敛并不排除途中越界\\
    \bottomrule
  \end{tabularx}
  \caption{同一数学工具进入不同问题时需要重新辨认的对象。投影、谱、固定点及轨迹
  条件的正式定义与证明已在前述章节建立；这里比较它们各自控制什么，不把相似公式
  当作可互换的保证。}
  \label{tab:future-reused-mathematical-objects}
\end{table*}

八类任务的共同起点是写清输入、未知对象和所求结论。选择函数后要控制总体误差，
规定表示后要检查算子关系，执行计算后要区分求解和统计误差；预测集合、扰动和隐私
再提出各自的量词要求。行动、来源和任务变化则改变证据如何产生、哪些旧结论仍可复用。
这里的回查引用指向通用定义和证明，当前任务所需的计算与条件保留在本章。

\section{评价数据选择后的总体误差}
\label{sec:future-learning-theory-generalization}

\subsection{候选选择与总体风险}
\label{sec:future-selection-population-risk}

学习理论要回答的不是“训练误差能否算出”，而是“用训练数据选出的函数在新样本上
会怎样”。选择动作使问题发生了变化：对一个事先固定的函数成立的概率界，
未必能直接用于看过同一批数据后才选出的函数。第~\ref{chap:probability-theory}章
建立的同时误差控制，正是从单个随机平均走向数据依赖选择的入口。

设非空有限候选类$\mathcal{H}=\{h_1,\ldots,h_M\}$在抽样前固定，样本
$S=(Z_1,\ldots,Z_n)$独立同分布，损失$\ell(h,Z)\in[0,1]$。
总体风险（population risk）是对未来同分布样本的平均损失，经验风险
（empirical risk）是手中样本的平均损失。记这两个量及各自的最优候选为
\begin{align*}
  R(h)&=\mathbb{E}[\ell(h,Z)],\\
  \widehat{R}_n(h)
  &=
  \frac{1}{n}\sum_{i=1}^{n}\ell(h,Z_i),\\
  \widehat{h}
  &\in
  \argmin_{h\in\mathcal{H}}\widehat{R}_n(h),\\
  h^\star
  &\in
  \argmin_{h\in\mathcal{H}}R(h).
\end{align*}
经验最优解依赖全部样本，不能在选择完成后把它当作预先固定的候选。定理
~\ref{thm:probability-finite-class-uniform-bound}已经从Hoeffding不等式和
联合界证明：以至少$1-\delta$的概率，全部$h\in\mathcal H$同时满足
\[
  |R(h)-\widehat R_n(h)|
  \leq
  \varepsilon_n(\delta)
  =
  \sqrt{\frac{\log(2M/\delta)}{2n}}.
\]
在这个同时事件上，把$\widehat h$从总体风险换到经验风险，使用经验最优性与
$h^\star$比较，再换回总体风险，便得到
\begin{equation}
  \begin{aligned}
    R(\widehat h)
    &\leq \widehat R_n(\widehat h)+\varepsilon_n\\
    &\leq \widehat R_n(h^\star)+\varepsilon_n
    \leq R(h^\star)+2\varepsilon_n(\delta).
  \end{aligned}
  \label{eq:future-erm-risk-bound}
\end{equation}
中间一步只用了经验最优性，两侧的转换才用了随机误差控制。误差事件必须
同时覆盖数据选出的$\widehat h$与用于比较的$h^\star$，不能在选完以后只证明其中一个。

实际求解器未必返回精确经验最优解。若$\widetilde h\in\mathcal H$只保证
\[
  \widehat{R}_n(\widetilde h)
  \leq
  \inf_{h\in\mathcal H}\widehat{R}_n(h)+\eta,
\]
则沿同一条比较路径可得
\[
  R(\widetilde h)
  \leq
  R(h^\star)+2\varepsilon_n(\delta)+\eta.
\]
这里的$\eta$来自优化过程，$\varepsilon_n(\delta)$来自有限样本；两项能够进入同一个
上界，不表示它们可由同一种论证减小。

\begin{example}[有限候选中的选择代价]
  \label{ex:future-finite-hypothesis-selection}
  假设$M=8$、$n=200$，并取$\delta=0.05$。代入上式可得
  \[
    \varepsilon_n(\delta)
    =
    \sqrt{\frac{\log 320}{400}}
    \approx 0.120.
  \]
  因而这个简单界只保证
  $R(\widehat h)\leq R(h^\star)+0.240$。若把$M$误写成$1$，
  便会忽略从八个候选中选择所付出的同时控制代价。数值较松不表示推导无效；
  它提示后续研究要利用函数类结构，不能只看候选数。
\end{example}

这个风险差还不是全部学习误差。若真实最优预测规则不在$\mathcal{H}$中，
$R(h^\star)$与所有可测规则中的最优风险之间存在\term{逼近误差}；
有限样本使经验量偏离总体量，产生\term{估计误差}；训练算法若只返回近似经验最优解，
还会产生\term{求解误差}。第~\ref{chap:combinatorics}章用打散与组合维度描述候选的
标记能力，第~\ref{chap:functional-analysis-and-approximation}章用函数空间和覆盖
描述逼近及距离分辨率；第~\ref{chap:probability-theory}章再将这些结构接到随机样本误差。
第~\ref{chap:optimization}章则分析后一项。三个来源可以相加进入一个上界，
却不能因为都叫“误差”而由同一个定理控制。

\subsection{函数类复杂度、稳定性与信息控制}
\label{sec:future-function-classes-algorithms}

当$\mathcal{H}$无限时，有限并集界不再直接适用。组合数学的生长函数计算样本上的
不同标记模式，泛函分析的覆盖数计算给定距离与精度下需要多少代表，概率论的
Rademacher型复杂度衡量函数类与随机符号相配合的程度。它们用不同方式把函数类结构
接到同时误差控制，不能仅因都衡量“复杂度”就互换；
算法稳定性改为比较替换一个样本前后的输出；信息论方法则控制算法输出携带了多少
关于训练样本的信息。这些是不同证明路径，不要求同时成立。用有限覆盖代替原函数类时，
覆盖所用的度量还必须能够控制当前损失差；参数空间中的欧氏距离小，本身不保证预测损失
也接近。

例如以输出距离$\sup_x|h(x)-g(x)|$作覆盖，且损失对预测值为$L$-Lipschitz，
每个$h$都有预先固定的代表$g$与它相距至多$r$，则
\[
  |R(h)-\widehat R_n(h)|
  \leq |R(g)-\widehat R_n(g)|+2Lr.
\]
两项$Lr$分别来自总体平均和样本平均。这样才能把有限代表上的概率界传回整个函数类；
若只知道参数数组相近而没有输出或损失的连续性控制，这一步无从成立。

\begin{example}[有限选择器的两类控制]
  \label{ex:future-stability-information-finite-selector}
  继续使用例~\ref{ex:future-finite-hypothesis-selection}中的
  \(M=8,n=200\)，先把训练算法\(A\)视为确定性算法，并把输出记为
  \(W=A(S)\in\mathcal H\)。若对任意只相差一个样本的\(S,S'\)和任意测试点\(z\)，
  都有
  \[
    \bigl|
      \ell(A(S),z)-\ell(A(S'),z)
    \bigr|
    \leq\beta,
  \]
  就可用留一替换把稳定性系数接到期望泛化差。具体地，令\(Z_i'\)是\(Z_i\)的
  独立同分布副本，并以\(S^{(i)}\)表示把\(S\)中的\(Z_i\)替换为\(Z_i'\)所得的数据集。
  交换\(Z_i\)与\(Z_i'\)不改变联合分布，因此
  \begin{align*}
    &\left|
      \mathbb{E}\!\left[
        R(A(S))-\widehat R_n(A(S))
      \right]
    \right|
    \\
    &\quad=
    \left|
      \frac1n\sum_{i=1}^n
      \mathbb{E}\!\left[
        \ell(A(S^{(i)}),Z_i)-\ell(A(S),Z_i)
      \right]
    \right|
    \leq\beta.
  \end{align*}
  第~\ref{sec:opt-stochastic-gradient-stability}节的成对更新分析可以用于推导稳定系数，
  上述替换论证则说明该系数为何控制期望泛化差。
  若\(A\)含内部随机性，可以固定同一个随机种子耦合两次运行并要求上述界对种子成立，
  或改用对算法随机性取期望的稳定性定义；不能省略这一随机性约定。
  这里真正需要核对的是算法对换样本是否敏感。未经正则化的经验风险最小化可能在
  两个经验风险几乎相同的候选之间切换；若它们在某个\(z\)上的损失分别为\(0\)和\(1\)，
  那么一次换样本就可能使上式只能取\(\beta=1\)。候选只有八个并不自动带来小稳定性。

  信息控制考察同一输出\(W\)携带了多少训练样本信息。因为\(W\)至多取八个值，
  \[
    I(W;S)\leq H(W)\leq\log 8.
  \]
  对\([0,1]\)值损失，第~\ref{chap:information-theory-and-statistical-geometry}章的
  换测度界因而给出
  \[
    \left|
      \mathbb{E}\!\left[
        R(W)-\widehat R_n(W)
      \right]
    \right|
    \leq
    \sqrt{\frac{I(W;S)}{2n}}
    \leq
    \sqrt{\frac{\log8}{400}}
    \approx0.072.
  \]
  这个数利用了输出范围有限，却没有证明换一个样本时输出接近。它控制的是期望泛化差，
  也不同于例~\ref{ex:future-finite-hypothesis-selection}中以高概率成立的超额风险界。
  稳定性与互信息可以分析同一个选择器，但条件和结论不能互换。
\end{example}

上述三种界不能直接按数值大小排序，因为所控制的量和概率口径不同。
表~\ref{tab:future-generalization-guarantees}固定同一有限选择器，把每个结论的输入
条件与输出量并列。互信息得到较小的$0.072$，并没有因此替代$0.240$的高概率
超额风险界；一个是平均泛化差，另一个还比较了类内最优规则。

\begin{table*}[htbp]
  \centering
  \begin{tabularx}{\linewidth}{lXX}
    \toprule
    数学路径 & 还需核对的条件 & 在当前例中控制的量\\
    \midrule
    有限候选同时控制 & 固定候选类、独立同分布、损失在$[0,1]$ & 至少$0.95$概率下，
    $R(\widehat h)-R(h^\star)\leq0.240$\\
    换样本稳定性 & 任意换一个样本后，任意测试损失变化至多$\beta$ & 期望泛化差的绝对值至多$\beta$\\
    输出互信息 & 输出至多取$8$个值及相应换测度条件 & 期望泛化差的绝对值至多$0.072$\\
    \bottomrule
  \end{tabularx}
  \caption{同一教学选择器的不同泛化保证，样本数固定为$200$。稳定性系数$\beta$
  需由算法另行证明，不能从候选数量小推出；期望结论也不能直接当作每次训练均成立的
  高概率结论。}
  \label{tab:future-generalization-guarantees}
\end{table*}

\subsection{样本复杂度与计算边界}
\label{sec:future-learning-computation-boundaries}

若希望类内超额风险至多为$\epsilon$，式~\eqref{eq:future-erm-risk-bound}给出
一个充分样本量条件：
\[
  n\geq \frac{2\log(2M/\delta)}{\epsilon^2}.
\]
取原例的$M=8,\delta=0.05$和$\epsilon=0.1$，向上取整得到$n=1154$。
这不是最少样本数，而是当前证明足以认证的预算。若求解器只保证经验目标差$\eta$，
还须有$\eta<\epsilon$，并将分母中的$\epsilon$换成$\epsilon-\eta$。
增加样本可以缩小随机误差，却不能消掉固定的求解误差。

观测不足的下界要比较难以区分的分布：若两种数据生成规律对现有观测几乎给出
相同分布，却要求不同的最佳答案，任何算法都难以同时正确；信息论距离可以定量表达
这种区分限度\scite{math-cover2006information}。计算边界则固定输入编码和答案要求，
问程序需要多少运算。若一个候选在一个样本上的损失可按单位成本求得，显式枚举
上述$M$个候选需要$Mn$次损失计算；当$M$随描述长度指数增长时，
样本界中的$\log M$增长缓慢，并不意味着枚举可行。这个枚举成本本身也不是
任何算法都需要指数时间的证明，困难性仍须按第~\ref{chap:combinatorics}章的归约论证。

式~\eqref{eq:future-erm-risk-bound}的边界很明确：损失有界、样本独立同分布、
候选类在看见样本前固定，而且使用同一分布定义经验风险和总体风险。
数据依赖地构造候选类、反复查看验证结果、面对相关序列或发生分布迁移时，
都要改用相应的条件化、序贯控制或换测度工具。
\BookPartRef{part:learning-theory}{第二卷“基础、理论和可信性”的“机器学习理论”部分}
按二分类、多分类和实值任务进一步规定可学习性；阅读那些结论时，需要分别核对
候选的表示、样本预算、返回解的质量和计算资源，而不是只比较某一个误差上界。

\section{按任务选择表示与算子结构}
\label{sec:future-representation-structure}

前节已经比较了数据选择与总体误差。现在先固定候选对象：有的任务从响应数据中拟合一个预测，
有的任务压缩数据结构，还有的任务规定输入变换应怎样传到输出。投影与核解决表示空间的选择，
局部算子、图关系与状态解决输入结构怎样写入函数；两类约束共同组成候选，而不自动形成学习保证。

\subsection{用投影、谱与核选择拟合和压缩的对象}
\label{sec:future-classical-kernel-geometric-representation}

上一节的泛化结论比较数据依赖的候选选择。要落实一次选择，还需固定函数族或表示目标。
本节沿投影、核空间和低维几何比较这些具体对象：同一矩阵可以是预测设计，也可以是
待压缩的数据，数学分解相同，研究目标与误差解释需要重新给出。

\subsubsection{投影拟合、谱结构与正则化}
\label{sec:future-fitting-projection-regularization}

经典模型经常共享平方损失、距离或谱分解，但共享运算不表示共享研究目标。
回归关心预测量与响应之间的误差，降维关心表示保留了什么结构，聚类则需要规定
何种群组关系才算成功。先看一个近共线回归例子，可以同时辨认拟合、敏感性与
正则化分别由哪一部分数学控制。

令设计矩阵与响应为
\[
  \bm{X}
  =
  \begin{bmatrix}
    1 & 1\\
    1 & 1+\varepsilon
  \end{bmatrix},
  \qquad
  \bm{y}
  =
  \begin{bmatrix}
    1\\
    1
  \end{bmatrix},
\]
其中$\varepsilon$很小。两列特征几乎相同，
因此某些系数组合对预测的作用也几乎相同。无正则最小二乘解满足正规方程
\[
  \bm{X}^{\top}\bm{X}\widehat{\bm{w}}
  =
  \bm{X}^{\top}\bm{y}.
\]
矩阵$\bm{X}^{\top}\bm{X}$的行列式为$\varepsilon^2$，
当$\varepsilon\to0$时趋近奇异。第~\ref{chap:linear-algebra-and-matrix-analysis}章
“线性空间”从两方面解释这个问题：投影确定最近的可达输出
$\bm{X}\widehat{\bm{w}}$，奇异值与条件数判断反求系数对微小扰动是否敏感。
两者分别对应第~\ref{sec:linear-algebra-reachability-solvability}节的求解任务与
第~\ref{sec:matrix-analysis-conditioning-perturbation}节的可靠性判断。

加入岭正则项后，目标变为
\[
  \min_{\bm{w}\in\mathbb{R}^{2}}
  \frac{1}{2}\lVert\bm{X}\bm{w}-\bm{y}\rVert_2^2
  +
  \frac{\lambda}{2}\lVert\bm{w}\rVert_2^2,
  \qquad \lambda>0,
\]
其唯一解为
\begin{equation}
  \widehat{\bm{w}}_\lambda
  =
  \left(
    \bm{X}^{\top}\bm{X}+\lambda\bm{I}
  \right)^{-1}
  \bm{X}^{\top}\bm{y}.
  \label{eq:future-ridge-solution}
\end{equation}
若$\sigma_j>0$是$\bm{X}$的奇异值，那么正规矩阵相应谱方向上的逆因子从
$1/\sigma_j^2$变为$1/(\sigma_j^2+\lambda)$；计入右端$\bm X^\top\bm y$后，
从响应到系数的放大因子则从$1/\sigma_j$变为$\sigma_j/(\sigma_j^2+\lambda)$。
正则化抑制了小奇异值方向的放大，但也改变了拟合目标。
它改善的是代数条件与系数稳定性，不自动证明新样本上的总体风险更小。

\begin{example}[共线特征的可辨方向]
  \label{ex:future-collinear-ridge}
  取$\varepsilon=0$时，两列完全相同，所有满足$w_1+w_2=1$的向量都给出
  $\bm{X}\bm{w}=\bm{y}$，所以预测值唯一而系数不唯一。加入$\lambda>0$后，
  严格凸性保证岭目标有唯一解，对称性使该解满足$w_1=w_2$，于是
  \[
    \widehat w_{\lambda,1}
    =
    \widehat w_{\lambda,2}
    =
    \frac{2}{4+\lambda}.
  \]
  此时
  \[
    \bm X\widehat{\bm w}_{\lambda}
    =
    \frac{4}{4+\lambda}\bm y,
  \]
  所以有限$\lambda$下的岭解不是零残差解，而是以收缩偏差换取较稳定的系数。
  零残差解集合中的最小范数解是$(1/2,1/2)^{\top}$；只有当
  $\lambda\downarrow0$时，岭解才趋于这个代表。这个结论仍不能说明两个特征具有
  相同的现实含义。
\end{example}

这里投影的对象是给定数据中的响应向量，而条件期望投影的对象是随机变量。
例如$X$为二值观测，$Y=X+\xi$，其中$\xi$独立于$X$、均值为零、方差为$1$。
在平方可积、只利用$X$的信息的预测中，最优者为$\mathbb E[Y\mid X]=X$，
其总体平方误差仍为$\mathbb E[\xi^2]=1$。正交性在此是
$\mathbb E[(Y-X)g(X)]=0$，对所有平方可积的$g(X)$成立；
它不是有限设计矩阵的等式$\bm X^\top(\bm y-\bm X\widehat{\bm w})=\bm0$。
两种投影都最小化平方距离，但前者使用分布下的内积，后者使用样本向量的内积。
即使某个有限样本能够零残差插值，也没有消去前者的不可预测噪声。

\subsubsection{有限特征与核空间计算}
\label{sec:future-finite-features-kernel-space}

同一个拟合思想可以从有限维系数扩展到函数空间。设$\mathcal{H}_K$是核$K$对应的
再生核Hilbert空间，取$\lambda>0$，考虑
\[
  \min_{f\in\mathcal{H}_K}
  \sum_{i=1}^{n}
  \bigl(f(\bm{x}_i)-y_i\bigr)^2
  +
  \lambda\lVert f\rVert_{\mathcal{H}_K}^{2}.
\]
把任意$f$正交分解为
\[
  f=f_{\parallel}+f_{\perp},
  \qquad
  f_{\parallel}
  \in
  \operatorname{span}
  \{K(\bm{x}_1,\cdot),\ldots,K(\bm{x}_n,\cdot)\}.
\]
再生性质给出
\[
  f_{\perp}(\bm{x}_i)
  =
  \langle
    f_{\perp},K(\bm{x}_i,\cdot)
  \rangle_{\mathcal{H}_K}
  =0.
\]
因此$f$与$f_{\parallel}$在训练点上的损失相同，而
\[
  \lVert f\rVert_{\mathcal{H}_K}^{2}
  =
  \lVert f_{\parallel}\rVert_{\mathcal{H}_K}^{2}
  +
  \lVert f_{\perp}\rVert_{\mathcal{H}_K}^{2}
  \geq
  \lVert f_{\parallel}\rVert_{\mathcal{H}_K}^{2}.
\]
任何最优解都可取成
\begin{equation}
  f^\star(\cdot)
  =
  \sum_{i=1}^{n}\alpha_i K(\bm{x}_i,\cdot).
  \label{eq:future-kernel-representation}
\end{equation}
这段投影推导解释了为什么无限维候选仍能化为$n$个系数，
而不是声称无限维问题本身已经消失。核矩阵的条件数、正则项和求解器仍决定
数值计算是否可靠。

令Gram矩阵$\bm K$的元素为$K(\bm x_i,\bm x_j)$，代入展开式后，
目标变为$\|\bm K\bm\alpha-\bm y\|_2^2+\lambda\bm\alpha^\top\bm K\bm\alpha$。
因为$\bm K$半正定，选择
\[
  (\bm K+\lambda\bm I)\bm\alpha=\bm y
\]
的解便满足这个二次目标的最优性条件；$\bm K+\lambda\bm I$正定，故该方程可解且解唯一。
即使$\bm K$奇异、同一函数存在多组展开系数，这个方程也给出一个确定代表。
例如$\bm K=\bigl[\begin{smallmatrix}1&1/2\\1/2&1\end{smallmatrix}\bigr]$、
$\lambda=1/2$、$\bm y=(1,0)^\top$时，
$\bm\alpha=(3/4,-1/4)^\top$，训练点预测为$\bm K\bm\alpha=(5/8,1/8)^\top$。
核空间中的正交分解因此落实为有限计算，但收缩后的预测仍不等于观测响应。
核与再生性质的主定义见第~\ref{sec:functional-positive-kernels-rkhs}节。

改变核会改变哪些函数被视为平滑或相近，改变正则项则改变同一函数空间中的偏好。
这两种改变分别作用于表示和选择准则。第~\ref{chap:differential-geometry-and-symmetry}章
把距离推广到局部度量、流形与对称结构，
第~\ref{chap:graph-theory}章把邻域关系写成图并用图谱分析变化模式。
它们可用于组织降维或聚类，却不能仅凭“距离保留”推出类别、语义或因果关系也被保留。

\subsubsection{聚类与低维几何表示}
\label{sec:future-prediction-clusters-low-dimensional-relations}

同一数据矩阵也可以作为无监督分析的对象，但这时目标已经改变。为使误差可以逐项复核，
在上例取\(\varepsilon=1\)，于是
\[
  \bm{X}=
  \begin{bmatrix}
    1&1\\
    1&2
  \end{bmatrix},
  \qquad
  \bm{X}^{\top}\bm{X}=
  \begin{bmatrix}
    2&3\\
    3&5
  \end{bmatrix}.
\]
\(\bm{X}^{\top}\bm{X}\)的两个特征值为
\[
  \sigma_1^2=\frac{7+3\sqrt5}{2},
  \qquad
  \sigma_2^2=\frac{7-3\sqrt5}{2}.
\]
若截断SVD只保留第一项
\(\bm{X}_1=\sigma_1\bm{u}_1\bm{v}_1^{\top}\)，
定理~\ref{thm:matrix-analysis-best-low-rank}给出
\begin{equation}
  \lVert\bm{X}-\bm{X}_1\rVert_{\mathrm F}^2
  =
  \sigma_2^2
  =
  \frac{7-3\sqrt5}{2}
  \approx0.146.
  \label{eq:future-truncated-svd-reconstruction-error}
\end{equation}
这个数是数据矩阵的秩一重构误差，不是回归误差。回归仍使用同一
未中心化矩阵\(\bm{X}\)和\(\bm{y}=(1,1)^{\top}\)，并最小化
\(\lVert\bm{X}\bm{w}-\bm{y}\rVert_2^2\)；此时
\(\bm{w}=(1,0)^{\top}\)使误差为零。若把两行
\((1,1)\)与\((1,2)\)视为样本做不预先中心化的一维降维，式
\eqref{eq:future-truncated-svd-reconstruction-error}衡量的是到所选秩一线性表示的
平方重构误差。若改做单簇聚类，目标又变为到质心
\(\bm m=(1,3/2)^{\top}\)的簇内平方和，其值为
\[
  \lVert(1,1)^{\top}-\bm m\rVert_2^2
  +
  \lVert(1,2)^{\top}-\bm m\rVert_2^2
  =
  \frac12.
\]
输入数据没有改变，但监督响应、低秩表示与离散分组三种输出接口不同，因而三个数不能
互相充当成功标准。

本节例子的结论止于给定样本、给定核与给定正则参数下的表示和求解。
噪声模型、预测区间、超参数选择及总体风险需要统计假设；
聚类和降维还须另行定义群组或结构保真的评价目标。
投影、谱和核是可复用工具，不是这些任务共同的成功标准。

\subsection{把输入关系写进局部与复合算子}
\label{sec:future-neural-structure-representation}

\subsubsection{局部共享运算与平移等变性}
\label{sec:future-local-operations-convolution-equivariance}

神经网络的结构研究关心怎样把输入中的先验关系写进函数族。
空间平移、节点重编号和时间推进都可表现为“变换输入后，输出应按可预测方式变化”，
但三者作用的对象不同。一个短序列上的共享局部算子可以把张量运算、卷积、
群作用和边界条件连接起来。

考虑长度为$4$的循环序列
\[
  \bm{x}
  =
  \begin{bmatrix}
    x_0 & x_1 & x_2 & x_3
  \end{bmatrix}^{\top}
\]
以及长度为$2$的核$\bm{k}=[k_0,k_1]^{\top}$。
沿用定义~\ref{def:signal-analysis-circular-convolution}的循环卷积，当前算子为
\[
  (T_{\bm{k}}\bm{x})_i
  =
  k_0x_i+k_1x_{(i-1)\bmod 4}.
\]
令$S$为循环右移一格的置换矩阵，即
$(S\bm{x})_i=x_{(i-1)\bmod4}$。直接代入可得
\begin{align}
  (T_{\bm{k}}S\bm{x})_i
  &=
  k_0x_{(i-1)\bmod4}
  +
  k_1x_{(i-2)\bmod4}
  \nonumber\\
  &=
  (ST_{\bm{k}}\bm{x})_i.
  \label{eq:future-convolution-equivariance}
\end{align}
所以$T_{\bm{k}}S=ST_{\bm{k}}$。
输入平移后，输出发生同样平移，这就是该算子对循环平移的\term{等变性}：
它不要求输出保持不变，而要求输出按规定的群作用同步变化。

\begin{example}[边界条件改变结构关系]
  \label{ex:future-convolution-boundary}
  取$\bm{x}=[1,2,0,0]^{\top}$、$\bm{k}=[1,-1]^{\top}$，
  则
  \[
    T_{\bm{k}}\bm{x}
    =
    [1,1,-2,0]^{\top}.
  \]
  循环右移输入得到$[0,1,2,0]^{\top}$，
  再卷积得到$[0,1,1,-2]^{\top}$，恰好是原输出的循环右移。
  若把循环边界改成左端补零，首尾不再属于同一平移轨道，
  上述交换关系在边界处失效。具体地，以$T_0$记补零的差分算子，改取输入
  $\bm e_3=(0,0,0,1)^\top$，则
  $T_0S\bm e_3=(1,-1,0,0)^\top$，而$ST_0\bm e_3=(1,0,0,0)^\top$。
  共享参数本身不足以保证任意边界处理下的精确等变。
\end{example}

第~\ref{chap:signal-analysis}章还说明，循环卷积在离散傅里叶坐标中化为逐频率乘法。
这同时给出表示和计算两种认识：每个频率模式被核的频率响应独立缩放，
而快速变换可以减少某些卷积的运算量。采样或池化随后会改变可分辨频率；
没有抗混叠条件时，低采样率中的低频模式可能来自原信号的多个高频成分。
结构约束限制了函数族，快速算法组织了计算，采样条件决定信息保留，
三者不能合并成一句“卷积保留结构”。

\subsubsection{图关系、重编号与状态表示}
\label{sec:future-sequences-graphs-state-representations}

输入若由无规则关系构成，变换不再是循环平移，而是节点重编号。
图上的映射$\Phi$若满足
\[
  \Phi(\bm{P}\bm{X},\bm{P}\bm{A}\bm{P}^{\top})
  =
  \bm{P}\Phi(\bm{X},\bm{A})
\]
对允许的置换矩阵$\bm{P}$成立，才具有相应的置换等变性。
第~\ref{chap:graph-theory}章提供邻接、路径、分离与图谱对象，
第~\ref{chap:differential-geometry-and-symmetry}章提供群作用下不变与等变的统一语言。
这条等式只说明重编号一致性，不说明映射能够区分所有不同图结构。

例如最简单的邻居求和为$\Phi(\bm X,\bm A)=\bm A\bm X$，重编号后
$(\bm P\bm A\bm P^\top)(\bm P\bm X)=\bm P\bm A\bm X$，直接使用了
$\bm P^\top\bm P=\bm I$。在三节点路径上取标量特征$(1,2,4)^\top$，
邻居和为$(2,5,2)^\top$。交换首尾节点后，特征变为$(4,2,1)^\top$，
输出仍为$(2,5,2)^\top$，恰是原输出的对应置换。
这个计算也暴露了限制：两个端点自身特征不同，邻居和却相同。
只用此层输出不能恢复被求和操作丢掉的信息，等变性并不是可区分性的证明。

持续到达的序列又引入状态。若
\[
  \bm{s}_{t+1}
  =
  \bm{A}\bm{s}_t+\bm{B}\bm{x}_t,
\]
展开后可得
\[
  \bm{s}_{t}
  =
  \bm{A}^{t}\bm{s}_0
  +
  \sum_{j=0}^{t-1}
  \bm{A}^{t-1-j}\bm{B}\bm{x}_j.
\]
第~\ref{chap:dynamical-systems-and-state-estimation}章的谱与Lyapunov工具
判断历史扰动如何衰减或放大。这个状态核与卷积形式相似，
但当$\bm{A}$随输入或时间改变时，固定卷积核不再描述完整递推。

状态转移中的谱还须与数据压缩中的谱区分。取两个转移矩阵
\[
  \bm A=\begin{bmatrix}1/2&4\\0&1/2\end{bmatrix},
  \qquad
  \bm C=\tfrac12\bm I,
  \qquad
  \bm s_0=(0,1)^\top,
\]
并令外部输入为零。二者的特征值都为$1/2$，但
\[
  \bm A^t\bm s_0=2^{-t}(8t,1)^\top,
  \qquad
  \bm C^t\bm s_0=2^{-t}(0,1)^\top.
\]
第一条轨迹在$t=1$时范数约为$4.03$，第二条已降为$0.5$。
二者最终都衰减，却只有后一条在欧氏范数下逐步缩小。
非正规矩阵的耦合产生短时放大；协方差矩阵是对称半正定的，
其正交谱分解并没有这种非正交模式耦合。图~\ref{fig:future-same-spectrum-transients}
把相同长期衰减率与不同短时响应分开。

\begin{figure}[htbp]
  \centering
  \begin{tikzpicture}[
  x=1mm,y=1mm,font=\figurefont\footnotesize,
  every node/.style={text=FigureInk,inner sep=1pt},
  axis/.style={draw=FigureStroke,line width=.5pt,->}
]
  \path[use as bounding box] (0,0) rectangle (112,65);
  \node[anchor=west] at (8,60) {$\|\bm s_t\|_2$};
  \begin{scope}[shift={(15,12)},x=10mm,y=10mm]
    \foreach \v in {1,2,3,4} {
      \draw[FigureGrid,densely dotted,line width=.4pt] (0,\v)--(8,\v);
      \node[left,xshift=-1mm] at (0,\v) {$\v$};
    }
    \draw[axis] (0,0)--(8.7,0) node[right] {$t$};
    \draw[axis] (0,0)--(0,4.6);
    \node[below left] at (0,0) {$0$};
    \foreach \t in {2,4,6,8}
      \node[below] at (\t,0) {$\t$};
    \draw[FigureRed,line width=.9pt]
      plot[domain=0:8,samples=9] (\x,{pow(.5,\x)*sqrt(1+64*\x*\x)});
    \draw[FigureBlue,dashed,line width=.9pt]
      plot[domain=0:8,samples=9] (\x,{pow(.5,\x)});
    \foreach \t in {1,...,8} {
      \fill[FigureRed] (\t,{pow(.5,\t)*sqrt(1+64*\t*\t)})
        circle[radius=.75mm];
      \path[draw=FigureBlue,fill=FigureBlue]
        (\t,{pow(.5,\t)}) ++(-.65mm,-.65mm)
        rectangle ++(1.3mm,1.3mm);
    }
    \path[draw=FigureStroke,fill=FigureYellow,line width=.4pt]
      (0,1) circle[radius=1mm];
    \node[anchor=west] at (3.3,3.5) {非正规 $\bm A$};
    \node[anchor=west] at (3.3,.7) {对角 $\bm C$};
  \end{scope}
\end{tikzpicture}

  \caption{同谱转移的不同瞬态。点为上述两条确定性轨迹在整数时刻的欧氏范数，
  线段仅连接离散时刻；初值、零输入和矩阵均由教学例指定，不是实验测量。
  相同特征值保证这里的最终衰减，却不保证相同的短时幅度。}
  \label{fig:future-same-spectrum-transients}
\end{figure}

\subsubsection{复合映射的敏感性与逼近能力}
\label{sec:future-composition-sensitivity-approximation}

结构确定后，还要判断复合映射对输入扰动的局部响应。
对
\[
  \bm{h}=f(\bm{x}),
  \qquad
  \bm{y}=g(\bm{h}),
\]
链式法则给出
\[
  \bm{J}_{g\circ f}(\bm{x})
  =
  \bm{J}_{g}(\bm{h})\bm{J}_{f}(\bm{x}),
\]
进而有局部范数界
\[
  \lVert\bm{J}_{g\circ f}(\bm{x})\rVert_2
  \leq
  \lVert\bm{J}_{g}(\bm{h})\rVert_2
  \lVert\bm{J}_{f}(\bm{x})\rVert_2.
\]
它控制局部敏感性，却不等价于第~\ref{chap:functional-analysis-and-approximation}章
讨论的逼近能力。一个函数族可能足以逼近目标，但具体参数对应的映射仍很敏感；
也可能满足严格等变性，却因结构约束过强而无法表达任务所需关系。

本节检查了输入关系如何经过运算保留或丢失。研究具体层、状态或邻域聚合时，
应先写出允许的输入变换，再检验交换等式、信息损失与计算代价。
等变性不推出不变性，局部Jacobian矩阵界不推出全局鲁棒性，
状态稳定也不推出表示包含了预测所需的全部信息。

\section{区分求导、求解收敛与预测泛化}
\label{sec:future-neural-training-generalization}

\subsection{反向求导、下降与曲率}
\label{sec:future-backpropagation-descent-curvature}

结构规定了候选函数怎样计算，训练则要从数据中选出参数。
这一步至少包含四个不同问题：导数是否算对，单步更新是否下降，
带噪迭代是否长期稳定，以及最终函数是否能泛化。
它们共享同一参数，却由不同数学条件支撑。

考虑一个标量输入、单隐藏单元的两层映射
\[
  f_{w,v}(x)=v\sigma(wx),
\]
以及单样本平方损失
\[
  L(w,v)
  =
  \frac{1}{2}\bigl(f_{w,v}(x)-y\bigr)^2.
\]
记误差$e=v\sigma(wx)-y$。在$\sigma$可微处，
链式法则给出
\begin{equation}
  \frac{\partial L}{\partial v}
  =
  e\sigma(wx),
  \qquad
  \frac{\partial L}{\partial w}
  =
  ev\sigma'(wx)x.
  \label{eq:future-two-layer-gradient}
\end{equation}
这只是导数计算。若$\sigma$在某点不可微，需要选用次梯度、
平滑替代或明确的人为替代梯度；三种做法改变的数学对象并不相同。

这个具体网络是否满足下降引理的光滑条件，可以直接由二阶导数核对。仍令
\(\bm\theta=(w,v)^{\top}\)，并把\(\sigma,\sigma',\sigma''\)都理解为在\(wx\)处取值。
在\(\sigma\)二阶可微处，
\[
  \nabla^2 L(w,v)
  =
  \begin{bmatrix}
    v^2x^2(\sigma')^2+evx^2\sigma''
    &
    x\sigma'(2v\sigma-y)
    \\
    x\sigma'(2v\sigma-y)
    &
    \sigma^2
  \end{bmatrix}.
\]
固定\(x,y\)，若某个包含当前参数的凸紧邻域\(U\)使\(wx\)始终落在
\(\sigma\)的二阶连续可微区间内，那么\(w,v,\sigma,\sigma'\)和\(\sigma''\)
在\(U\)上都有界，以上Hessian矩阵的算子范数也有有限上界\(L_U\)。
第~\ref{chap:mathematical-analysis}章的中值估计于是给出
\[
  \lVert\nabla L(\bm\theta)-\nabla L(\bm\theta')\rVert_2
  \leq
  L_U\lVert\bm\theta-\bm\theta'\rVert_2,
  \qquad
  \bm\theta,\bm\theta'\in U.
\]
所以这个两层网络在该邻域内满足梯度Lipschitz条件，但要使用局部常数\(L_U\)，
还须保证更新没有离开\(U\)。即使\(\sigma,\sigma',\sigma''\)全局有界，
上式中的\(v^2\)也说明不限制\(v\)时不能由此得到全局光滑常数；若使用ReLU且
\(x\neq0\)，跨过\(wx=0\)的邻域中一般不存在上述经典Hessian矩阵。

设把参数合并为$\bm\theta=(w,v)^{\top}$，
目标$F(\bm\theta)$的梯度是$L_F$-Lipschitz的。
梯度更新
\[
  \bm{\theta}^{(t+1)}
  =
  \bm{\theta}^{(t)}
  -
  \eta\nabla F(\bm{\theta}^{(t)})
\]
结合光滑性上界可得
\begin{align}
  F(\bm{\theta}^{(t+1)})
  &\leq
  F(\bm{\theta}^{(t)})
  -
  \eta\lVert\nabla F(\bm{\theta}^{(t)})\rVert_2^2
  +
  \frac{L_F\eta^2}{2}
  \lVert\nabla F(\bm{\theta}^{(t)})\rVert_2^2
  \nonumber\\
  &=
  F(\bm{\theta}^{(t)})
  -
  \eta
  \left(1-\frac{L_F\eta}{2}\right)
  \lVert\nabla F(\bm{\theta}^{(t)})\rVert_2^2.
  \label{eq:future-gradient-descent-one-step}
\end{align}
当$0<\eta<2/L_F$时，右端不增。
第~\ref{chap:optimization}章给出了这种下降论证及凸性、PL条件等进一步接口；
非凸目标仅有单步下降，通常不足以推出到达全局最优解。

\begin{example}[步长跨过下降范围]
  \label{ex:future-training-step-size}
  对$F(\theta)=\theta^2/2$，梯度为$\theta$且$L_F=1$。
  更新后$\theta^{(t+1)}=(1-\eta)\theta^{(t)}$，
  因而
  \[
    \frac{F(\theta^{(t+1)})}{F(\theta^{(t)})}
    =
    (1-\eta)^2.
  \]
  取$\eta=1$时一步到达零点，取$\eta=3$时目标放大为原来的四倍。
  梯度方向没有算错，失败来自步长不满足下降条件。
\end{example}

曲率矩阵的谱决定不同方向上的尺度差异。若局部二次模型中的Hessian矩阵
\(\bm H=\bm H^{\top}\succ\bm0\)，且近似为
\[
  F(\bm{\theta}+\bm{\Delta})
  \approx
  F(\bm{\theta})
  +
  \langle\nabla F(\bm{\theta}),\bm{\Delta}\rangle
  +
  \frac{1}{2}
  \bm{\Delta}^{\top}\bm{H}\bm{\Delta},
\]
那么\(\bm H\)的最大与最小特征值之比较刻画各方向的曲率尺度；条件数较大时，
统一步长受最大曲率方向限制，而在小曲率方向上推进较慢。取对称正定矩阵
\(\bm P\in\mathbb{R}^{2\times2}\)，预条件更新
\[
  \bm{\theta}^{(t+1)}
  =
  \bm{\theta}^{(t)}
  -
  \eta\bm{P}^{-1}\nabla F(\bm{\theta}^{(t)})
\]
试图改变局部坐标中的尺度。第~\ref{chap:matrix-analysis}章说明
$\bm{P}$的谱与数值求解条件；具体下降保证仍须结合预条件后的谱界和步长。
若\(\bm H\)不定，则还存在负曲率方向，不能沿用正定二次模型的条件数解释，
下降与收敛也需要另行分析。

\subsection{小批量随机逼近与分布几何}
\label{sec:future-minibatch-stochastic-approximation-geometry}

第~\ref{chap:information-theory-and-statistical-geometry}章把Fisher度量解释为参数
变化引起的局部分布变化。
经验梯度的二阶矩只有在相应概率模型与取期望条件下才与Fisher信息对应，
不能仅因矩阵外形相同便互换。

\begin{example}[同一参数点上的三种二次量]
  \label{ex:future-parameter-second-moment-fisher}
  沿用上一小节的两层映射，取\(\sigma(u)=u\)、\(x=1\)，并把
  \(f_{w,v}(x)=vw\)作为Bernoulli模型的对数几率：
  \[
    p_{\bm\theta}(Y=1\mid x)
    =
    \frac{1}{1+\exp\{-f_{w,v}(x)\}}.
  \]
  在\(\bm\theta_0=(1,\log4)^{\top}\)处，正类概率为\(p=0.8\)，且
  \[
    \bm J
    =
    \nabla_{\bm\theta}f_{\bm\theta_0}(x)
    =
    \begin{bmatrix}\log4\\1\end{bmatrix}.
  \]
  对同一小扰动\(\bm\Delta=(\varepsilon,0)^{\top}\)，欧氏参数距离只记录坐标变化，
  \[
    \lVert\bm\Delta\rVert_2^2=\varepsilon^2.
  \]

  若当前观察到\(Y=1\)，负对数似然梯度为
  \(\bm g_1=(p-1)\bm J=-0.2\bm J\)。这个观测形成的梯度二阶矩为
  \(\bm G_{\mathrm{obs}}=\bm g_1\bm g_1^{\top}\)，故
  \[
    \bm\Delta^{\top}\bm G_{\mathrm{obs}}\bm\Delta
    =
    0.04(\log4)^2\varepsilon^2
    \approx0.0769\varepsilon^2.
  \]
  Fisher信息则对模型自身的随机标签取期望：
  \[
    \bm F(\bm\theta_0)
    =
    \mathbb{E}_{Y\sim p_{\bm\theta_0}}
    [\bm g_Y\bm g_Y^{\top}]
    =
    p(1-p)\bm J\bm J^{\top},
  \]
  因而
  \[
    \bm\Delta^{\top}\bm F(\bm\theta_0)\bm\Delta
    =
    0.16(\log4)^2\varepsilon^2
    \approx0.3075\varepsilon^2.
  \]
  三个量在同一参数点和同一方向上仍不相等：欧氏距离依赖参数坐标，
  观测梯度二阶矩依赖实际样本，Fisher二次型则是模型分布下得分的期望，
  它的一半给出局部KL变化的二阶主项。只有损失确为该模型的负对数似然，并按模型分布
  对观测取期望时，后两种矩阵才对应。
\end{example}

小批量训练把精确梯度改成随机估计
\[
  \bm{g}^{(t)}
  =
  \nabla F(\bm{\theta}^{(t)})
  +
  \bm{\xi}^{(t)}.
\]
这里$\mathcal F_t$包含第$t$次抽取小批量之前的信息。若固定训练集，从中独立、
有放回地抽取$b$个样本梯度，并且单次梯度对经验目标$F$条件无偏、
条件噪声二阶矩至多为$\sigma^2$，则平均后的噪声满足
\[
  \mathbb E[\bm\xi^{(t)}\mid\mathcal F_t]=\bm0,
  \qquad
  \mathbb E[\|\bm\xi^{(t)}\|_2^2\mid\mathcal F_t]\leq \frac{\sigma^2}{b}.
\]
这个$1/b$来自条件独立性；相关采样不能直接沿用。将随机梯度代入下降引理，
在$F$全局$L_F$-光滑、$0<\eta\leq1/L_F$时得到
\[
  \begin{aligned}
    \mathbb E[F(\bm\theta^{(t+1)})\mid\mathcal F_t]
    &\leq F(\bm\theta^{(t)})
    -\frac{\eta}{2}\|\nabla F(\bm\theta^{(t)})\|_2^2\\
    &\quad+\frac{L_F\eta^2\sigma^2}{2b}.
  \end{aligned}
\]
若初始参数固定且$F\geq F_{\min}$，累加$T$步便有
\[
  \frac1T\sum_{t=0}^{T-1}\mathbb E\|\nabla F(\bm\theta^{(t)})\|_2^2
  \leq
  \frac{2(F(\bm\theta^{(0)})-F_{\min})}{\eta T}
  +\frac{L_F\eta\sigma^2}{b}.
\]
固定步长留下噪声项，增大批量可缩小该项，却没有改变目标是经验风险这一事实。
若采用递减步长，还需核对步长求和、矩条件及迭代稳定性；
第~\ref{sec:stoch-approximation-markov-noise}节说明随机逼近如何与平均动力学联系，
以及相关观测为何需要额外条件。

\subsection{训练收敛与泛化条件}
\label{sec:future-training-convergence-generalization}

即使训练目标持续下降，也只说明当前数据上的求解行为。
把训练所得函数用于新样本，需要回到
式~\eqref{eq:probability-finite-class-bound}一类泛化工具，
或采用函数类复杂度、稳定性和信息控制等其他路径。
参数数量、参数范数、谱量和优化轨迹可以成为分析量，
但任何一个量单独变小都不自动推出总体风险变小。

一个极端反例能排除“收敛所以泛化”的推断。令$X$在$[0,1]$均匀分布，
$Y=1$恒成立，采用零一损失。算法记住有限训练输入，在这些点输出$1$，
在其余点输出$0$，以后不再更新。它的训练损失为零，迭代输出也已经固定；
新输入落入有限训练点集的概率却为零，所以总体风险等于$1$。
这里没有违反有限类风险界：候选必须允许任意有限记忆集合，
不能把训练结束后那个单一输出当作抽样前固定的候选类。
这个反例不是在描述通常的神经网络训练，而是在指出仅有求解与收敛信息不足以作统计判断。

研究具体训练方法时，初始化、数据变换和停止规则会共同决定输出的分布；
它们必须进入函数类或算法分析。
本节的边界是：式~\eqref{eq:future-two-layer-gradient}只保证局部导数正确，
式~\eqref{eq:future-gradient-descent-one-step}只在光滑性和步长条件下保证单步变化，
随机逼近讨论长期迭代，而泛化仍需单独的数据生成假设。

\section{区分条件查询、参数学习与样本生成}
\label{sec:future-probabilistic-inference-learning-generation}

神经网络训练的例子主要把参数作为优化变量。给定联合概率模型后，
还要区分三个对象：固定参数下的条件分布、由数据推断的参数，以及生成的新随机样本。
概率与图论支持第一个计算，统计与优化支持第二个过程，随机过程与动力学支持第三个过程；
三个过程共享模型却没有相同的误差保证。

\subsection{联合分布、条件查询与结构计算}
\label{sec:future-joint-model-conditional-query}

概率模型中的“计算概率”“从数据估计参数”和“生成新样本”常使用同一个联合分布，
却回答三个不同问题。先固定一个含隐变量的简单模型，
可以看清边缘化、条件化、参数学习和采样各自改变了什么。

设隐变量$Z\in\{0,1\}$，观测$X\in\{0,1\}$，
\[
  \mathbb{P}(Z=1)=\pi,
  \qquad
  \mathbb{P}(X=1\mid Z=z)=\theta_z.
\]
参数$\pi,\theta_0,\theta_1$给定时，边缘概率由全概率公式得到
\[
  \mathbb{P}(X=1)
  =
  (1-\pi)\theta_0+\pi\theta_1.
\]
以下条件查询要求$(1-\pi)\theta_0+\pi\theta_1>0$。
若该量为零，$X=1$在当前模型下是不可能事件，不能用概率比值唯一确定其后验。
在这一正概率条件下，观察到$X=1$后，对隐状态的条件推断为
\begin{equation}
  \mathbb{P}(Z=1\mid X=1)
  =
  \frac{\pi\theta_1}
  {(1-\pi)\theta_0+\pi\theta_1}.
  \label{eq:future-latent-posterior}
\end{equation}
这里没有学习参数，只是对给定联合分布执行求和与条件化。

\begin{example}[同一联合分布上的三个问题]
  \label{ex:future-latent-three-questions}
  取$\pi=1/4$、$\theta_0=1/5$、$\theta_1=4/5$。
  则
  \[
    \mathbb{P}(X=1)
    =
    \frac{3}{4}\frac{1}{5}
    +
    \frac{1}{4}\frac{4}{5}
    =
    \frac{7}{20},
  \]
  且
  \[
    \mathbb{P}(Z=1\mid X=1)
    =
    \frac{(1/4)(4/5)}{7/20}
    =
    \frac{4}{7}.
  \]
  推断问的是已见$X=1$时$Z$如何分布；学习会把三个参数视为未知；
  生成则先抽$Z$，再按对应$\theta_Z$抽$X$。三个过程共用因子，
  输入、输出和正确性标准却不同。
\end{example}

变量增多后，直接枚举联合状态的代价可能迅速增大。
第~\ref{chap:graph-theory}章中的分离、消元与分配结构
把全局求和重组为局部运算。计算量取决于消元过程中形成的中间作用域，
而不只取决于原图有多少条边。连续变量还要核对积分交换与可积性；
线性高斯结构可利用第~\ref{chap:matrix-analysis}章的Schur补，
动态结构则可复用第~\ref{chap:dynamical-systems-and-state-estimation}章的
预测、滤波与平滑递推。

同一联合分解还会因查询目标不同而使用不同运算。计算某个变量的边缘概率时，
需要对其余变量求和；求联合最大后验配置时，要把消元中的求和改为最大化，并保存
每一步的最优见证以便回溯出彼此兼容的联合配置。逐时边缘最优状态则先分别计算各时刻的
边缘分布，再在每个边缘上取最大值；这些状态通常不能直接拼成联合概率最大的整条路径。
求和与最大化一般不可交换，因此图结构只规定可利用的分解，查询目标还决定消元代数、
回溯信息和最终输出。

\subsection{变分、采样与局部曲率近似}
\label{sec:future-exact-query-three-approximations}

精确边缘难以计算时，可以引入近似分布$q(z)$。
假设$0<p(x)<\infty$，并且$q$与后验$p(\cdot\mid x)$具有相同支撑，即
$q(z)>0$当且仅当$p(x,z)>0$；再假设下述期望与KL散度有限，则
\begin{align}
  \log p(x)
  &=
  \log
  \sum_z
  q(z)\frac{p(x,z)}{q(z)}
  \nonumber\\
  &\geq
  \sum_z
  q(z)
  \log\frac{p(x,z)}{q(z)}
  =
  \mathcal{L}(q),
  \label{eq:future-elbo}
\end{align}
其中不等式来自对数函数的凹性。进一步整理可得
\[
  \log p(x)-\mathcal{L}(q)
  =
  \operatorname{KL}\bigl(q(z)\,\|\,p(z\mid x)\bigr).
\]
因此优化变分下界等价于在所选$q$族内逼近后验，
但只有当$q$族包含真实后验且优化到全局最优时，间隙才必为零。
第~\ref{chap:information-theory-and-statistical-geometry}章提供
KL散度与变分表达，第~\ref{chap:optimization}章提供交替和分布优化的求解接口。

另一条路径是用样本平均近似期望。若
$Z^{(1)},\ldots,Z^{(N)}$来自目标分布且$f(Z)$的二阶矩有限，
\[
  \widehat\mu_N
  =
  \frac{1}{N}
  \sum_{j=1}^{N}f(Z^{(j)})
\]
在独立抽样下具有方差
$\operatorname{Var}(f(Z))/N$。
若样本来自Markov链，相关性会改变有效方差，
需要第~\ref{sec:stoch-markov-long-run-average}节的遍历与相关平均条件。
链保持目标分布并不说明有限步已经充分混合，
有限样本误差也不能由平稳性一项保证。

连续后验还可以在众数附近作局部二次展开。若
\[
  \widehat{\bm z}
  \in
  \argmax_{\bm z}\log p(\bm z\mid x),
  \qquad
  \bm H
  =
  -\nabla^2_{\bm z}\log p(\widehat{\bm z}\mid x)
\]
且$\widehat{\bm z}$是内点、$\bm H$正定，二阶Taylor展开给出局部高斯近似
\[
  p(\bm z\mid x)
  \approx
  \mathcal N(\widehat{\bm z},\bm H^{-1}).
\]
第~\ref{sec:statistics-model-evidence}节的Laplace近似说明，还需有足够光滑性，并且
后验主要质量确由该邻域控制。多峰、边界众数或退化曲率都会破坏这一近似。
变分法限制分布族，采样法以随机平均逼近期望，Laplace方法只保留局部曲率；
三者改变的对象和误差来源不同。

\begin{example}[同一后验上的三种近似]
  \label{ex:future-posterior-approximation-comparison}
  取先验\(p\sim\operatorname{Beta}(2,2)\)，观察四次Bernoulli试验中的三次成功，
  则同一个精确后验为
  \[
    p\mid\bm y\sim\operatorname{Beta}(5,3),
    \qquad
    \mathbb{E}[p\mid\bm y]=\frac58,
    \qquad
    \operatorname{Var}(p\mid\bm y)=\frac{15}{576}.
  \]
  以下记其密度为\(\pi(p\mid\bm y)\)。
  若变分族取全部Beta分布，ELBO的全局最优解就是
  \(q^\star=\operatorname{Beta}(5,3)\)，其分布逼近误差
  \(\operatorname{KL}(q^\star\Vert \pi(\cdot\mid\bm y))\)为零。这个结果依赖候选族
  包含精确后验；换成受限族后，ELBO间隙衡量的是所选方向的KL误差。

  若能独立采样，\(N\)个后验样本的均值估计后验均值，其Monte Carlo标准误为
  \[
    \sqrt{\frac{15}{576N}}.
  \]
  取\(N=100\)时约为\(0.0161\)。这里控制的是某个后验期望的随机积分误差，
  不是一个固定近似密度与精确后验之间的距离。

  Laplace方法在后验众数\(\widehat p=2/3\)处展开。负对数后验的局部曲率为
  \[
    -\left.
    \frac{\mathrm d^2}{\mathrm dp^2}
    \log \pi(p\mid\bm y)
    \right|_{p=2/3}
    =
    \frac{4}{(2/3)^2}+\frac{2}{(1/3)^2}
    =27,
  \]
  因而得到\(\mathcal N(2/3,1/27)\)。其均值相对精确后验偏高
  \(1/24\)，方差相差
  \[
    \frac1{27}-\frac{15}{576}
    =
    \frac{19}{1728}
    \approx0.0110.
  \]
  它还把少量质量放到参数区间\([0,1]\)之外。Laplace误差来自局部二次形状无法完整
  表示偏斜与边界，而非有限样本平均的随机波动。
\end{example}

表~\ref{tab:future-posterior-approximations}保留同一个精确后验作为比较基准。
变分误差回用第~\ref{chap:information-theory-and-statistical-geometry}章的分布距离；
采样误差与局部高斯近似则同属第~\ref{chap:random-variables}章“随机变量”的推断工具，
分别回查第~\ref{sec:probability-monte-carlo-variance-reduction}节的随机积分和
第~\ref{sec:statistics-model-evidence}节的Laplace近似。
增加采样数会减少随机积分误差，却不会自动改善一个受限变分族或局部高斯形状。

\begin{table*}[htbp]
  \centering
  \begin{tabularx}{\linewidth}{lXX}
    \toprule
    近似路径 & 当前构造的结果 & 误差含义\\
    \midrule
    全部Beta分布的变分族 & 全局最优为$\operatorname{Beta}(5,3)$ & 分布KL为零；族包含真后验且优化精确\\
    $100$个独立后验样本 & 样本均值近似$5/8$ & Monte Carlo标准误约$0.0161$，估计一个期望\\
    众数附近二次展开 & $\mathcal N(2/3,1/27)$ & 局部形状误差；均值偏高$1/24$且越出参数区间\\
    \bottomrule
  \end{tabularx}
  \caption{教学后验$\operatorname{Beta}(5,3)$上的三种近似。采样一行没有给出实际抽取
  的样本值，标准误由精确后验方差计算；表中数字不是实验结果，也不是三种误差同量纲
  的性能排名。}
  \label{tab:future-posterior-approximations}
\end{table*}
\FloatBarrier

\subsection{参数学习与生成分布}
\label{sec:future-parameter-estimation-sample-generation}

参数未知时，问题转为根据观测识别和估计
$\pi,\theta_0,\theta_1$。仅观察$X$时，
不同参数组合可能产生同一个边缘Bernoulli概率，
所以更多样本只能更精确地估计这个边缘量，不能自动区分不可识别的参数。
第~\ref{sec:rv-inference}节把识别、似然、缺失机制和估计误差分开；
交替优化只能求解给定目标，不能修复目标本身的不可识别性。

在这个隐变量模型中，交替精确更新为何单调改善观测似然也可以就地核对。对观测
\(x_1,\ldots,x_n\)，记参数为\(\bm\vartheta=(\pi,\theta_0,\theta_1)\)，并定义
\[
  \ell_{\mathrm{obs}}(\bm\vartheta)
  =
  \sum_{i=1}^{n}
  \log\sum_{z\in\{0,1\}}
  p_{\bm\vartheta}(x_i,z).
\]
设各次迭代中的边缘似然严格为正且有限，后验和下列期望均有定义，并且两个条件
子问题的最优解能够达到。E步固定参数，计算每个样本来自各隐状态的后验权重；
M步固定这些权重，重新选择参数以提高加权完整数据目标。精确更新分别取
\[
  q_i^{(t+1)}(z)
  =
  p_{\bm\vartheta^{(t)}}(z\mid x_i),
  \qquad
  \bm\vartheta^{(t+1)}
  \in
  \argmax_{\bm\vartheta}
  \mathcal L(\bm q^{(t+1)},\bm\vartheta),
\]
其中
\[
  \mathcal L(\bm q,\bm\vartheta)
  =
  \sum_{i=1}^{n}\sum_z
  q_i(z)
  \log
  \frac{p_{\bm\vartheta}(x_i,z)}{q_i(z)}.
\]
变分恒等式给出
\[
  \ell_{\mathrm{obs}}(\bm\vartheta)
  =
  \mathcal L(\bm q,\bm\vartheta)
  +
  \sum_{i=1}^{n}
  \operatorname{KL}\!\left(
    q_i\,\middle\|\,
    p_{\bm\vartheta}(\,\cdot\mid x_i)
  \right).
\]
E步使右端的KL项在\(\bm\vartheta^{(t)}\)处为零，M步又精确提高同一个下界，因而
\begin{align*}
  \ell_{\mathrm{obs}}(\bm\vartheta^{(t)})
  &=
  \mathcal L(\bm q^{(t+1)},\bm\vartheta^{(t)})
  \\
  &\leq
  \mathcal L(\bm q^{(t+1)},\bm\vartheta^{(t+1)})
  \\
  &\leq
  \ell_{\mathrm{obs}}(\bm\vartheta^{(t+1)}).
\end{align*}
这就是单调改善的完整理由；一般交替与主化--最小化论证见
第~\ref{chap:optimization}章。等号可以连续成立，所以该链条不保证严格改善，
也不保证到达全局最大值。若E步不是当前参数的精确后验，或M步没有提高同一个
\(\mathcal L\)，上述两处依据都要重新检查。

沿用例~\ref{ex:future-latent-three-questions}的初始参数，若只有两次独立观测
$x_1=1,x_2=0$，E步给出的$Z=1$权重分别为$4/7$与$1/13$。
M步按加权频率更新：
\begin{align*}
  \pi^{\mathrm{new}}&=\frac{4/7+1/13}{2}=\frac{59}{182},\\
  \theta_1^{\mathrm{new}}&=\frac{4/7}{4/7+1/13}=\frac{52}{59},\\
  \theta_0^{\mathrm{new}}&=\frac{3/7}{3/7+12/13}=\frac{13}{41}.
\end{align*}
新参数的边缘成功概率为$1/2$，两次观测的似然从
$(7/20)(13/20)=91/400$提高到$1/4$。但
$\theta_0=\theta_1=1/2$与任意$\pi$也给出这个边缘分布；
似然提高并没有告诉我们哪一组隐状态解释是真实的。

生成样本又改变了输出要求。可逆变换依赖Jacobian矩阵与密度换元，
连续随机生成依赖布朗运动、随机微分与密度演化，
而直接祖先采样依赖图中的条件分解。例如$U$在$[0,1]$均匀分布，
$X=1+2U$的密度在$[1,3]$上为$1/2$，来自逆变换导数
$|\mathrm du/\mathrm dx|=1/2$；忽略这个因子，所得函数积分便不是$1$。
对前面的二值模型，先按$\pi$抽$Z$再按$\theta_Z$抽$X$，才会产生相应边缘成功率；
直接把后验概率$4/7$当作生成$X$的概率，则混淆了给定观测后的查询和模型自身的采样。

一个过程具有正确的不变分布、
一个变分目标足够高，或一个数值积分器局部误差较小，
都不能单独替代对生成分布与任务质量的评价。
若模型边缘概率设错，增加采样数只会更准确地复现错误模型。
研究时因此要分别检查联合模型是否合适、条件查询近似是否准确、
参数是否可识别、随机样本是否服从所求分布，以及离散数值过程是否引入额外偏差。

\section{为覆盖、扰动与数据保护分别建立保证}
\label{sec:future-trustworthy-machine-learning}

\subsection{预测集合覆盖与验证证据}
\label{sec:future-prediction-error-coverage-evidence}

平均预测误差只回答随机抽取一个样本时的平均表现。
可信机器学习会进一步询问：预测范围是否覆盖未来结果，输入受扰动后输出是否稳定，
不同群组或干预目标是否得到规定保障，算法输出是否泄露单个数据主体的信息。
这些要求没有共同的单一分数，必须先固定保护或保证的对象。

先看有限样本预测覆盖。设有$n$个校准分数
$S_1,\ldots,S_n$和一个未来分数$S_{n+1}$，
且$n+1$个分数可交换。若分数越大表示预测越不相容，
令
\[
  k=\left\lceil(n+1)(1-\alpha)\right\rceil,
\]
当$k\leq n$时取排序阈值$q=S_{(k)}$，并令预测集合包含所有分数不超过$q$的候选结果；
当$k=n+1$时约定$q=+\infty$，即返回整个候选结果空间。后一约定处理
$\alpha<1/(n+1)$的极端置信水平，此时有限校准样本无法给出非平凡的分布无关集合。
由于未来分数的秩在无并列的情形下均匀分布于
$\{1,\ldots,n+1\}$，
\begin{equation}
  \mathbb{P}(S_{n+1}\leq q)
  =
  \begin{cases}
    k/(n+1), & k\leq n,\\
    1, & k=n+1,
  \end{cases}
  \geq
  1-\alpha.
  \label{eq:future-exchangeable-coverage}
\end{equation}
有并列时采用适当随机化或保守规则，仍可得到相应覆盖结论。

\begin{example}[四个可交换分数的秩]
  \label{ex:future-conformal-rank}
  设$n=3$、$\alpha=1/4$，三个校准分数排序后为
  $0.2,0.5,0.9$。此时$k=\lceil4\times3/4\rceil=3$，
  阈值为$0.9$。若四个分数可交换且没有并列，
  未来分数落在前三级秩中的概率为$3/4$。
  这个计算给出边缘覆盖；它没有声称对每个固定输入、每个群组或每个时间点
  都具有$3/4$的条件覆盖。
\end{example}

式~\eqref{eq:future-exchangeable-coverage}依赖可交换性与预先规定的构造。
若校准样本与未来样本来自不同分布，或者反复查看校准表现后选择分数函数，
秩的对称性会被破坏。第~\ref{chap:random-variables}章“随机变量”把这些条件放在
同一条推断链上：第~\ref{sec:probability-exchangeability-conformal-ranks}节由可交换秩
建立预测覆盖，第~\ref{sec:stats-uncertainty-testing-selection}节再说明模型选择、
多重比较与持续查看怎样影响证据。总体覆盖不能仅通过按群组报告平均值
升级为条件覆盖。

同样，例~\ref{ex:future-conformal-rank}中的$3/4$是在校准样本与未来样本一起重复
抽取时成立的边缘概率。即使当前已经看到阈值$0.9$，秩对称性也没有单独确定“给定这
一次校准结果”的未来条件覆盖。平均训练误差、边缘预测覆盖和固定群组条件覆盖使用
不同的随机对象及量词，不能只因数值都写成一个概率就互换。


\subsection{扰动鲁棒性、群组与行为约束}
\label{sec:future-perturbations-groups-behavior}

鲁棒性把问题改成允许输入在集合内变化时，输出最多改变多少。
若可微标量函数$f$在$\bm{x}$附近作一阶展开，
\[
  f(\bm{x}+\bm{\delta})-f(\bm{x})
  \approx
  \langle\nabla f(\bm{x}),\bm{\delta}\rangle.
\]
对约束$\lVert\bm{\delta}\rVert_p\leq\varepsilon$，
对偶范数关系给出
\begin{equation}
  \sup_{\lVert\bm{\delta}\rVert_p\leq\varepsilon}
  \langle\nabla f(\bm{x}),\bm{\delta}\rangle
  =
  \varepsilon
  \lVert\nabla f(\bm{x})\rVert_q,
  \qquad
  \frac{1}{p}+\frac{1}{q}=1.
  \label{eq:future-dual-norm-robustness}
\end{equation}
第~\ref{chap:linear-algebra}章给出范数与对偶范数，
第~\ref{chap:optimization}章把扰动集合接入约束或鞍点目标。
式~\eqref{eq:future-dual-norm-robustness}是线性项的精确最坏变化；
对有限扰动下的非线性模型，还要控制Taylor余项或给出有效的全局Lipschitz界。

例如线性得分$f(\bm x)=x_1+2x_2$在$\ell_\infty$半径$0.1$内的最大变化为
$0.1\|(1,2)^\top\|_1=0.3$。若当前得分为$0.5$、分类阈值为零，
整个扰动集合内仍为正类；若得分只有$0.2$，同一半径就没有这个保证。
对非线性$f$，若它在整个扰动球上连续可微且
$\sup_{\bm z}\|\nabla f(\bm z)\|_q\leq L$，沿线段积分才能推出有限变化
$|f(\bm x+\bm\delta)-f(\bm x)|\leq L\varepsilon$。
只在中心点算一个导数，不足以完成这一步；预测集合的边缘覆盖也没有控制这个最坏扰动。

公平性可能要求比较群组条件错误率，也可能要求某种干预或反事实量。
前者依赖条件概率及分母，后者还依赖第~\ref{chap:causal-inference}章的识别假设。
仅让两个观测群组的某项平均相等，不能推出个体在干预意义下受到相同处理。
行动补救也不只是寻找梯度方向：候选改变必须可实施，
并且改变后预测的因果含义需要由结构假设支持。

群组约束也不能从总体均值中读出。若群组$A$占$90\%$且错误率为零，
群组$B$占$10\%$且错误率为一，总体错误率只有$10\%$，却完全没有保护$B$。
研究群组差异时必须报告相应条件总体、分母与不确定性；
把样本很少的群组估计值直接当成真实错误率，又会引入统计证据不足的问题。

对齐、拒答和安全决策把概率输出继续接到损失、行动与风险。
第~\ref{chap:decision-theory-and-dynamic-risk}章说明同一预测分布可以对应
不同错误代价和风险准则，第~\ref{chap:control-theory}章则区分长期稳定、
成本最优与安全集合保持。

例如二分类错误成本为$10$、正确成本为零、拒答成本为$0.3$，
在真实条件正类概率$p$下，三种行动的条件风险分别为
$10(1-p)$、$10p$和$0.3$。当$p=0.6$时应拒答，当$p=0.98$时预测正类的风险为$0.2$，
低于拒答成本。这个结论来自明确的损失表，不是从“置信较高”直接推出；
若输入的$p$只是失准的模型分数，相应行动也未必风险最小。

\subsection{差分隐私、解释与干预对象}
\label{sec:future-data-changes-privacy-explanation}

隐私保护比较指定数据变化下随机算法的输出分布。需要先约定保护单位：
例如固定记录数时，只替换一条记录的两个数据集称为相邻，记为$D\sim D'$。
这是记录级保护；若一个人贡献多条记录，而目标是保护这个人，就应重新规定相邻关系，
比较该人全部记录的增删或替换。这里“相邻”是允许的数据变化，不是数据数值彼此接近。
其专门保证见
\BookDefinitionRef{def:privacy-dp}{第二卷“基础、理论和可信性”的“隐私保护”章“差分隐私”定义}。
这里复用事件概率的比较。给定$\epsilon\geq0$、$0\leq\delta<1$，对随机机制$\mathcal M$，
需要核对每一对相邻$D,D'$和每个输出事件$E$是否满足
\[
  \mathbb P(\mathcal M(D)\in E)
  \leq e^\epsilon\mathbb P(\mathcal M(D')\in E)+\delta.
\]
概率来自机制的内部随机性；相邻关系对称时，交换$D,D'$后的不等式也必须成立。
若某个机制只输出$0$或$1$，在一对
相邻数据集上的输出$1$概率分别为$0.4$和$0.6$，则两个单点事件在两个方向上的
概率比都不超过$1.5$。因而针对这一对数据集，乘性界可取$e^\epsilon=1.5$、
$\delta=0$；要获得机制的隐私保证，还须对全部相邻数据集核验，不能只检查这一对。
一般的$\delta$是事件概率比较中的加性裕量，不能直接理解为某个人被完全公开的概率。

相邻关系也决定一个查询最多会改变多少。对标量查询$q$，其敏感度定义为
\[
 \Delta q=\sup_{D\sim D'}|q(D)-q(D')|.
\]
例如每条记录$z_i\in[0,1]$，查询$q(D)=\sum_i z_i$在单条替换下的敏感度为$1$；
若同一人的记录数至多为$m$，并将相邻改为增删该人的全部记录，最大变化可达$m$。
对输出$q(D)+W$这类给查询加随机噪声的常见机制，满足同一隐私要求所需的噪声尺度
必须按相应最大变化校准；一般差分隐私机制不必采用这种加噪形式。
研究者还须说明保护谁、攻击者能看到哪些输出及辅助信息，这些约定构成这里的威胁模型。
分布差异和信息控制可以帮助证明给定模型下的保证，却不能代替这些约定。
类似地，可解释性中的局部导数、估计影响和联盟贡献分别解释
输入敏感性、训练过程与预先定义的价值函数，不能共用一个含混的“重要性”概念。

\begin{example}[同一预测上的三种解释对象]
  \label{ex:future-explanation-three-objects}
  考虑带交互项的预测
  \[
    f_{\bm w}(\bm x)
    =
    w_1x_1+w_2x_2+w_{12}x_1x_2,
    \qquad
    \bm w=(1,2,1)^{\top},
  \]
  并解释目标输入\(\bm x=(1,1)^{\top}\)处的预测\(f_{\bm w}(\bm x)=4\)。
  局部导数把训练结果\(\bm w\)固定，只改变输入的第一个特征：
  \[
    \frac{\partial f_{\bm w}}{\partial x_1}(1,1)
    =
    w_1+w_{12}x_2
    =
    2.
  \]
  它表示\(x_1\)在该点每作一个无穷小单位变化时，预测的一阶变化率。

  影响函数改变的是训练分布。为使数值可复核，设单点损失为半平方损失
  \[
    \ell(\bm w;(\bm x,y))
    =
    \frac12\bigl(f_{\bm w}(\bm x)-y\bigr)^2,
  \]
  \(\bm w\)由相应总体风险估计，且该总体风险的参数Hessian矩阵在当前解处为
  \(\bm I_3\)。若给训练点
  \(z=((1,0)^{\top},4)\)增加无穷小概率质量，则该点的参数梯度为
  \[
    \nabla_{\bm w}\ell(\bm w;z)
    =
    \bigl(f_{\bm w}(1,0)-4\bigr)
    \begin{bmatrix}1\\0\\0\end{bmatrix}
    =
    \begin{bmatrix}-3\\0\\0\end{bmatrix}.
  \]
  第~\ref{chap:mathematical-statistics}章的影响函数公式给出参数的一阶变化
  \(3(1,0,0)^{\top}\)，进而使目标输入处的预测以速率\(3\)变化。
  这里被扰动的是训练点权重，而不是目标输入的\(x_1\)。

  Shapley解释还要先指定缺失特征的基线。取零基线并令
  \(v(S)=f_{\bm w}(\bm x_S,\bm0_{\bar S})-f_{\bm w}(\bm0)\)，则
  \[
    v(\varnothing)=0,\quad
    v(\{1\})=1,\quad
    v(\{2\})=2,\quad
    v(\{1,2\})=4.
  \]
  第~\ref{chap:game-theory-and-multi-agent-decision}章的两人Shapley公式因此给出
  \[
    \phi_1
    =
    \frac12[(1-0)+(4-2)]
    =
    \frac32,
    \qquad
    \phi_2=\frac52.
  \]
  它把相对基线的有限预测差\(4\)分配给两个特征，并把交互项在两种加入次序间平均。
  数值\(2,3,3/2\)不同并非矛盾：三者分别回答局部输入扰动、训练分布污染和
  基于指定联盟价值的有限贡献。
\end{example}

以上三种解释也都没有单独识别干预效果。将特征改成基线、对模型输入求导、
或重加权训练样本，不等于在真实机制中实施$\operatorname{do}(X_1=x)$。
要解释行动会造成什么后果，仍须指定结构因果模型和识别条件。
可信性研究应逐项规定保护对象、威胁条件、数据协议和验证程序；
覆盖、扰动和隐私推导不能被拼成一个无定义的总体可信保证。

\section{核对行动改变数据后的规划与学习条件}
\label{sec:future-reinforcement-learning-sequential-decision}

\subsection{两阶段价值、固定点与占用流量}
\label{sec:future-two-step-fixed-point-occupancy}

监督预测把数据分布视为给定，序贯决策中的行动却会改变后续状态和未来数据。
因此，价值递推、随机逼近、探索证据、离线支持和安全约束必须在同一个时间结构中
分别核对\scite{sutton2018}。先从一个已知模型的两阶段过程出发，可以确定动态规划
究竟解决了哪一步。

设初始状态为$s$，可选行动$a$，
即时奖励为$r(s,a)$，下一状态按
$P(s'\mid s,a)$分布。最后阶段的价值为
\[
  V_1(s)
  =
  \max_a r_1(s,a).
\]
初始阶段的最优价值则为
\begin{equation}
  V_0(s)
  =
  \max_a
  \left\{
    r_0(s,a)
    +
    \sum_{s'}P(s'\mid s,a)V_1(s')
  \right\}.
  \label{eq:future-two-stage-bellman}
\end{equation}
这一步来自条件期望和“在下一状态到达后再选择”的信息顺序。
若初始行动必须在看见$s'$之前同时承诺第二步行动，
最大值与期望的次序就会改变。

\begin{example}[即时奖励与未来价值]
  \label{ex:future-two-stage-decision}
  初始状态有行动$a$和$b$。行动$a$立即得到$1$，
  随后以相同概率到达终局价值$0$或$4$的状态；
  行动$b$立即得到$2$，随后确定到达终局价值$1$的状态。
  由式~\eqref{eq:future-two-stage-bellman}，
  \[
    Q_0(s,a)=1+\frac{0+4}{2}=3,
    \qquad
    Q_0(s,b)=2+1=3.
  \]
  两个行动的期望价值相同，但尾部风险不同。
  若改用动态风险准则，还要明确条件风险怎样逐阶段嵌套，
  不能只把期望符号替换为一个静态风险数。
\end{example}

沿用第~\ref{sec:decision-dynamic-programming}节的有限折扣问题，把成本改成奖励并同步将最小化改成最大化，Bellman最优算子写成
\[
  (\mathcal{T}V)(s)
  =
  \max_a
  \left\{
    r(s,a)
    +
    \gamma
    \sum_{s'}P(s'\mid s,a)V(s')
  \right\},
  \qquad 0\leq\gamma<1.
\]
成本与奖励的换号不改变第~\ref{sec:decision-dynamic-programming}节已证明的压缩系数，因此有
\[
  \lVert\mathcal{T}V-\mathcal{T}W\rVert_\infty
  \leq
  \gamma\lVert V-W\rVert_\infty.
\]
结合式~\eqref{eq:decision-value-iteration-error}，这个奖励约定同样得到固定点唯一性
与值迭代误差收缩。这个结论以已知转移核、奖励有界和$\gamma<1$为前提；
平均回报、无折扣或连续状态问题需要不同的空间与条件。

同一个有限状态、有限动作的折扣问题还可从流量角度表达。给定归一化初始分布
$\rho$，以$(1-\gamma)\sum_s\rho(s)V(s)$为最小化目标，并加入价值不等式
\[
  V(s)
  \geq
  r(s,a)
  +
  \gamma\sum_{s'}P(s'\mid s,a)V(s')
\]
组成一个线性规划；其对偶变量$d(s,a)\geq0$满足
\begin{equation}
  \sum_a d(s,a)
  =
  (1-\gamma)\rho(s)
  +
  \gamma
  \sum_{\bar s,\bar a}
  d(\bar s,\bar a)P(s\mid\bar s,\bar a).
  \label{eq:future-discounted-occupancy-flow}
\end{equation}
对所有状态求和可核对$\sum_{s,a}d(s,a)=1$，所以$d$表示归一化折扣
占用流量；对偶目标是最大化$\sum_{s,a}d(s,a)r(s,a)$。
第~\ref{chap:optimization}章的线性规划对偶与
第~\ref{chap:decision-theory-and-dynamic-risk}章的占用度量共同说明，
价值函数和访问流量是同一已知模型规划问题的两侧。转移核未知时，这个等价关系
并没有说明如何取得满足守恒约束所需的数据。

\subsection{随机观测下的价值与策略更新}
\label{sec:future-random-observations-value-policy-updates}

模型未知时，期望只能由交互数据近似。例如一次时序差分更新可写为
\[
  V^{(t+1)}(S_t)
  =
  V^{(t)}(S_t)
  +
  \alpha_t
  \left(
    R_t+\gamma V^{(t)}(S_{t+1})-V^{(t)}(S_t)
  \right).
\]
括号中的时序差分残差包含平均更新方向和随机误差，不能把它整体视为独立同分布噪声；
状态本身也由策略驱动的Markov链产生。
第~\ref{chap:probability-theory}章提供条件均值与鞅差语言，
第~\ref{sec:stoch-approximation-markov-noise}节把Markov噪声、步长条件与平均ODE联系起来，
并检查漂移稳定性。
同策略线性逼近下的稳定分析不能直接搬到离策略、非线性或持续变化的环境。

把这句话落实到一个可检查的平均漂移。令
\(V_{\bm\theta}(s)=\bm\phi(s)^{\top}\bm\theta\)，把各状态特征按行组成
\(\bm\Phi\)。固定策略\(\pi\)诱导有限状态转移矩阵\(\bm P_\pi\)，并假设样本链
已经处于平稳分布\(\bm d\)，记\(\bm D=\operatorname{diag}(\bm d)\)。
行为策略与被评价策略同为\(\pi\)时，
\[
  \overline{\bm h}(\bm\theta)
  =
  \mathbb{E}\!\left[
    \bm\phi(S_t)
    \bigl(
      R_t+\gamma\bm\phi(S_{t+1})^{\top}\bm\theta
      -\bm\phi(S_t)^{\top}\bm\theta
    \bigr)
  \right]
  =
  \bm b-\bm A_{\mathrm{TD}}\bm\theta,
\]
\[
  \bm b=\bm\Phi^{\top}\bm D\bm r_\pi,
  \qquad
  \bm A_{\mathrm{TD}}
  =
  \bm\Phi^{\top}\bm D
  (\bm I-\gamma\bm P_\pi)\bm\Phi.
\]
这里\(\bm r_\pi\)是固定策略在各状态下的一步期望奖励向量。
平稳性使\(\bm P_\pi\)在\(\bm D\)加权二范数下非扩张。若
\(0\leq\gamma<1\)且特征在实际访问分布下非退化，即
\(\bm\Phi^{\top}\bm D\bm\Phi\succ\bm0\)，那么对任意
\(\bm c\neq\bm0\)，
\[
  \bm c^{\top}\bm A_{\mathrm{TD}}\bm c
  \geq
  (1-\gamma)
  \lVert\bm\Phi\bm c\rVert_{\bm D}^{2}
  >
  0.
\]
因此\(\bm A_{\mathrm{TD}}\)的对称部分正定，平均ODE
\(\dot{\bm\theta}=\bm b-\bm A_{\mathrm{TD}}\bm\theta\)具有唯一的全局指数稳定平衡。
加权非扩张与Lyapunov核对已在
第~\ref{chap:dynamical-systems-and-state-estimation}章完整给出，这里不重复证明。
从平均ODE走回随机TD迭代，还需要
\(\sum_t\alpha_t=\infty\)、\(\sum_t\alpha_t^2<\infty\)，以及相应的噪声矩、
Markov混合和迭代有界条件。若没有平稳采样，有限时间均值还含初始分布的过渡项；
若改成离策略采样，\(\bm D\)与\(\bm P_\pi\)不再配套，上述非扩张步骤可能失效。

若目标改为直接调整参数化策略$\pi_{\bm\theta}$，取不依赖
$\bm\theta$的初始分布$\rho_0$和环境转移核$P$。有限时域轨迹
$\tau=(S_0,A_0,\ldots,A_{T-1},S_T)$的概率分解为
\[
  p_{\bm\theta}(\tau)
  =
  \rho_0(S_0)
  \prod_{t=0}^{T-1}
  \pi_{\bm\theta}(A_t\mid S_t)
  P(S_{t+1}\mid S_t,A_t).
\]
再令
\[
  G(\tau)
  =
  \sum_{t=0}^{T-1}\gamma^t r_t(S_t,A_t),
  \qquad
  J(\bm\theta)
  =
  \mathbb{E}_{\tau\sim p_{\bm\theta}}[G(\tau)].
\]
这里奖励函数$r_t$不显含$\bm\theta$，所以$G$对参数的依赖只来自轨迹分布。若
$G$可积、策略关于$\bm\theta$可微、策略支持集不随$\bm\theta$改变，并且允许交换
求导与积分，对上述乘积分布使用对数导数便得到
\begin{equation}
  \nabla_{\bm\theta}J(\bm\theta)
  =
  \mathbb{E}_{\tau\sim p_{\bm\theta}}
  \left[
    G(\tau)
    \sum_{t=0}^{T-1}
    \nabla_{\bm\theta}
    \log\pi_{\bm\theta}(A_t\mid S_t)
  \right].
  \label{eq:future-policy-gradient}
\end{equation}
若初始分布、环境动力学或奖励本身也显式依赖$\bm\theta$，相应对数导数或
\(\nabla_{\bm\theta}G_{\bm\theta}(\tau)\)项必须保留，不能只写策略得分之和。
若基线$b_t(S_t)$只依赖当前状态，不依赖尚未抽取的$A_t$，则
\[
  \mathbb{E}\!\left[
    b_t(S_t)\nabla_{\bm\theta}
    \log\pi_{\bm\theta}(A_t\mid S_t)
    \,\middle|\,S_t
  \right]
  =
  b_t(S_t)\nabla_{\bm\theta}\sum_a\pi_{\bm\theta}(a\mid S_t)
  =
  \bm0.
\]
因此这类基线可以降低方差而不改变期望梯度；依赖当前动作的任意基线没有这一保证。
第~\ref{chap:information-theory-and-statistical-geometry}章的Fisher度量进一步按
局部策略分布变化缩放式~\eqref{eq:future-policy-gradient}，所得自然梯度不能与
任意经验梯度二阶矩预条件混为一谈。

\subsection{自适应探索与学习复杂度}
\label{sec:future-adaptive-sampling-exploration-complexity}

探索又改变了数据生成策略。自归一化集中可以利用实际访问形成的设计矩阵，
但仍需可预测性、噪声条件和足够激励；信息论下界说明某些环境在有限交互中难以区分。

\begin{example}[自适应两动作探索中的置信与访问预算]
  \label{ex:future-two-action-exploration-budget}
  设单状态中有两个动作\(A,B\)，Bernoulli奖励的均值分别为
  \(\mu_A=0.6\)与\(\mu_B=0.5\)，并假设给定过去与当前动作后，奖励仍服从
  相应参数的Bernoulli分布。在第\(t\)轮，算法根据
  \(\mathcal F_{t-1}\)中的历史选择动作\(A_t\)，再观察奖励\(Y_t\)。
  对动作\(a\)，
  \[
    X_{t,a}
    =
    \mathbb{I}\{A_t=a\}(Y_t-\mu_a)
  \]
  是有界鞅差，其可预测方差累计量为
  \(\mu_a(1-\mu_a)N_a(t)\)，其中\(N_a(t)\)是截至\(t\)的实际访问次数。
  因而第~\ref{chap:stochastic-processes-and-approximation}章的Freedman界按
  \(N_a(t)\)及实际条件方差控制误差，而不是把随机访问序列误当作预先固定样本。

  两个动作还可写成一个自归一化线性模型。令
  \(\bm\theta^\star=(\mu_A,\mu_B)^{\top}\)，选择动作时取
  \(\bm x_t=\bm e_{A_t}\)，并写成
  \[
    Y_t
    =
    \langle\bm x_t,\bm\theta^\star\rangle+\eta_t.
  \]
  Bernoulli中心化噪声条件\(1/2\)-次高斯，且\(\bm x_t\)在当前奖励出现前已经确定。
  取正则参数\(\lambda=1\)和
  \(\lVert\bm\theta^\star\rVert_2\leq B=1\)，设计矩阵为
  \[
    \bm V_t
    =
    \bm I+
    \begin{bmatrix}
      N_A(t)&0\\
      0&N_B(t)
    \end{bmatrix}.
  \]
  令
  \[
    \widehat{\bm\theta}_t
    =
    \bm V_t^{-1}
    \sum_{s=1}^{t}\bm x_sY_s
    =
    \begin{bmatrix}
      \widehat\mu_A\\
      \widehat\mu_B
    \end{bmatrix}
  \]
  为相应的岭估计。自归一化置信椭球以至少\(1-\delta\)的概率对所有\(t\)同时满足
  \[
    \lVert
      \widehat{\bm\theta}_t-\bm\theta^\star
    \rVert_{\bm V_t}
    \leq
    \beta_t(\delta)
    =
    \frac12
    \sqrt{
      2\log
      \frac{
        \sqrt{(1+N_A(t))(1+N_B(t))}
      }{\delta}
    }
    +1.
  \]
  于是均值差\(\Delta=\mu_A-\mu_B\)的\emph{统计误差}由实际访问数决定：
  \[
    \left|
      (\widehat\mu_A-\widehat\mu_B)-\Delta
    \right|
    \leq
    \beta_t(\delta)
    \sqrt{
      \frac{1}{1+N_A(t)}
      +
      \frac{1}{1+N_B(t)}
    }.
  \]
  记上式右端为\(r_t\)，并记估计差为
  \(\widehat\Delta_t=\widehat\mu_A-\widehat\mu_B\)。
  取\(\delta=0.05\)。若某一时刻实际有
  \(N_A=N_B=2200\)，则\(\beta_t\approx3.312\)，且\(r_t\approx0.0998\)，
  略小于真实间隔\(\Delta=0.1\)。因此在该置信事件上
  \(\widehat\Delta_t\geq\Delta-r_t>0\)，点估计的符号保持正确。
  但投影置信区间\([\widehat\Delta_t-r_t,\widehat\Delta_t+r_t]\)是否排除零，
  还要检查实际观测是否满足\(\widehat\Delta_t>r_t\)。若只根据真实间隔和误差界
  预先保证整个置信集合都区分两个动作，一个充分条件是\(2r_t<\Delta\)；
  当前访问数并不满足这个更强条件。

  \emph{访问预算}是另一项要求。一个强制两边各访问\(2200\)次的协议需要至少
  \(4400\)次交互；仅给出总预算\(t=4400\)却不规定探索分配，可能出现
  \(N_B(t)=0\)，此时上述半径仍含\(1/(1+N_B(t))=1\)，无法辨认动作\(B\)。
  前一个不等式回答“给定已经实现的设计，统计误差多大”，探索协议则回答
  “用多少次交互才能实现足够的设计”。二者必须分开报告。
\end{example}

函数类复杂度或低秩结构要转化为样本效率结论，还须核对候选类包含什么真对象、
数据怎样获得，以及误差按什么残差衡量。
例如在有限状态价值学习中，可实现性可以指真实最优价值$V^\star$属于候选类$\mathcal H$；
若不属于，就还要计入逼近误差。这不是说每个行动都可执行，也不是说每个状态都已被访问。
不同学习问题可能要求真实价值、动作价值或环境模型属于相应候选类，必须分别指明对象。
对Bellman残差，先把候选价值与其Bellman更新之差定义为残差，
再在某个数据生成规则下取期望；若这个由“候选函数—数据规则”索引的二元量可分解为
$d$维向量内积，$d$才构成Bellman秩分析的入口。普通矩阵低秩本身既没有定义这个
残差，也不保证学习者会访问辨认它所需的状态与行动。

\subsection{信息、反馈与奖励识别}
\label{sec:future-information-feedback-identifiability}

部分观测要求把历史压缩为足以决策的状态或条件分布。
固定日志下的离线评价还要求目标策略的轨迹分布被日志策略覆盖；
若某个行动在日志中概率为零，重要性权重无法恢复其结果。
因果部分识别可以描述支持不足时仍可确定的范围，却不能凭借更复杂的拟合补出
从未观测的行动后果。

奖励并不总是被直接记录。给定轨迹回报
\[
  R_{\bm\theta}(\tau)
  =
  \sum_{t=0}^{T-1}\gamma^t
  r_{\bm\theta}(S_t,A_t),
\]
示范数据需要先规定奖励怎样影响被记录的轨迹分布。
令$\beta>0$表示回报敏感度，也称逆温度（inverse temperature）：
它越大，模型越偏向高回报；它对应
第~\ref{sec:information-gibbs-optimization}节所用温度的倒数。
对整段示范资料，可以考虑按回报对一个固定参考轨迹分布加权的Gibbs似然。
参考分布$\mu$须同时指定初态、环境转移及动作参考规则，不能只由环境动力学确定。
离散轨迹下，当$0<Z(\bm\theta,\beta)<\infty$时，定义
\[
  p_{\bm\theta,\beta}(\tau)
  =
  \frac{
    \mu(\tau)\exp\{\beta R_{\bm\theta}(\tau)\}
  }{
    Z(\bm\theta,\beta)
  }.
\]
这里$Z(\bm\theta,\beta)=\sum_\tau\mu(\tau)\exp\{\beta R_{\bm\theta}(\tau)\}$
对全部允许轨迹求和；连续轨迹时改为相对于参考测度的积分。
这一分布是对整段资料的候选统计模型。它可用于描述偏向高回报的记录选择，
但是否符合示范的实际生成或筛选方式，仍需另行核对。

合法的整轨迹分布不自动等于固定随机环境中某项在线策略的分布。
考虑两阶段模型，各状态只有一个行动：从初态以各$1/2$概率到达$g$或$b$，
初始奖励为$0$，第二阶段奖励分别为$2$和$0$，取$\gamma=1/2$。
两条轨迹的回报分别为$1$和$0$。以真实轨迹分布为$\mu$，在$\beta=1$时，上式给出
\[
 p_{\bm\theta,1}(\tau_g)=\frac{(1/2)e}{(1/2)e+1/2}
 =\frac{e}{1+e}\approx0.731.
\]
然而没有任何可选行动能把环境到达$g$的概率从$1/2$改成$0.731$。
指数加权同时偏向了环境碰巧给出好结果的轨迹。
因此若目标是解释固定环境中可执行的策略行为，应对满足前述
“初态、策略、环境转移”乘积分解的行为模型建立似然；
不能仅凭整轨迹归一化就认定其满足这些约束。

偏好反馈则可把同一回报差接到二元选择概率：
\begin{equation}
  \mathbb{P}_{\bm\theta,\beta}
  (\tau^+\succ\tau^-)
  =
  \frac{1}{
    1+\exp\{-\beta[
      R_{\bm\theta}(\tau^+)-R_{\bm\theta}(\tau^-)
    ]\}
  }.
  \label{eq:future-reward-preference-likelihood}
\end{equation}
例如回报差为$1$时，$\beta$从$1$增至$2$，
较高回报轨迹被偏好的概率从约$0.731$增至$0.881$，与逆温度的含义一致。
在给定参数后各条反馈条件独立的模型下，示范轨迹概率和成对选择概率的乘积构成
关于\((\bm\theta,\beta)\)的似然。
这把第~\ref{chap:mathematical-statistics}章的似然接口接到了奖励学习：数据固定后，
比较的是不同奖励参数对已观察行为赋予的概率。

似然能够优化并不等于奖励可识别。若两个参数对所有允许查询产生相同的示范与偏好
分布，那么只有在它们也给出相同目标量时，该目标才可识别。例如给所有轨迹回报加同一
常数不会改变归一化示范概率或偏好差，因而绝对奖励零点不可识别；若\(\beta\)也未知，
同时缩放回报并反向缩放\(\beta\)仍给出相同观测分布，奖励尺度也不可识别。只在已访问
轨迹上收集偏好还可能留下其他行动的回报差未定。研究目标因此应明确写成可识别的
回报差、行为分布、最优策略等价类或另加规范化与覆盖条件，而不能把任意一个拟合出的
\(r_{\bm\theta}\)当作唯一真实奖励。

\subsection{多方决策、动态风险与安全约束}
\label{sec:future-tasks-participants-risk-safety}

多方参与时，环境转移不再是固定背景。
第~\ref{chap:game-theory-and-multi-agent-decision}章中的最佳响应、均衡、信息集和反事实遗憾分别规定
允许怎样偏离以及对手如何进入比较。零和平均策略保证、一般和均衡存在与
某个学习动态的最后迭代收敛不是同一结论。
行为约束还会调用第~\ref{chap:decision-theory-and-dynamic-risk}章的动态风险、
第~\ref{chap:dynamical-systems-and-state-estimation}章的Lyapunov稳定性和
第~\ref{chap:control-theory}章的安全不变性。

固定点在这里也换了对象。折扣Bellman算子把一个实值函数映到另一个实值函数，
压缩性给出唯一解和迭代速率。Nash均衡要求每方策略属于对他方策略的最佳响应集合；
有限博弈中的混合策略单纯形虽紧凸，最佳响应却通常多值，存在性论证不等于压缩迭代。
例如两方分别偏好硬币相同与不同的零和博弈，唯一混合均衡为各以$1/2$选择正面；
若同步采用对上一轮纯行动的最佳响应，从$(1,1)$开始会循环经过
$(1,0),(0,0),(0,1),(1,1)$。均衡存在没有让最后一次策略收敛。

扩展式博弈的反事实遗憾则在给定信息集比较“在这里改选行动”的收益差，
按对手与机会到达该处的概率加权，并排除己方到达概率。
它是策略偏离的记账方式，仍须遵守玩家不可区分历史的信息约束；
它并非在结构因果模型中把某个变量方程替换为常量的静态干预。
有限、完美回忆等条件下的遗憾分解也不能脱离博弈类型，直接升级为一般和博弈的
最后迭代收敛保证。

动态风险需要保持同样清晰的信息顺序。例~\ref{ex:future-two-stage-decision}中，
行动$a$的总奖励为$1$或$5$，行动$b$的总奖励恒为$3$。
若把成本设为奖励的负值，采用各阶段条件最坏成本，即对当前信息下仍可能发生的
后继成本取最大值，则$a$的最坏成本为$-1$，$b$为$-3$，最小化成本会选择$b$。
在更长的树上，这种准则应递推为
$C_0+\rho_0(C_1+\rho_1(C_2))$，其中$\rho_t$只使用第$t$阶段可用的信息。
这里的逐节点条件最坏值支持向后递推；任意静态风险函数未必具有同样的时间一致性。
期望相同只回答平均回报相同，并未选择风险准则。

安全约束还会排除成本较低却越界的轨迹。回到共同温度模型，$x_t$为相对于环境的温差，
$u_t$为加热量。在无噪声时取
\begin{align*}
  x_{t+1}&=0.8x_t+u_t,\qquad r=2,\\
  u_t&=0.4-0.3(x_t-2)=1-0.3x_t.
\end{align*}
平衡输入为$0.4$，误差$e_t=x_t-2$满足$e_{t+1}=0.5e_t$。
若规定$x_t\in[0,3]$和$u_t\in[0,1]$，则对安全集合内每个状态，
$u_t\in[0.1,1]$，且$x_{t+1}=0.5x_t+1\in[1,2.5]$；
这才是一步安全保持的核验，而非只看误差收缩。
若再加入扰动$w_{t+1}$且逐步保证$|w_{t+1}|\leq0.5$，下一状态仍在$[0.5,3]$内，
可以归纳得到全过程安全。高斯噪声具有无界支撑，不能使用这一确定性保证，
只能在另行规定的概率约束或失效预算下讨论安全。

这些算例要求研究者先写清四件事：
优化什么回报或风险，数据由哪一个策略产生，决策时可以使用什么信息，
当前结论控制价值误差、遗憾、稳定性还是逐轨迹约束。

\section{判断来源证据怎样进入监督与决策}
\label{sec:future-knowledge-acquisition-fusion}

\subsection{来源可靠性、相关性与可识别性}
\label{sec:future-source-answers-reliability-identifiability}

训练标签、规则和偏好并不总是直接可得。
当多个来源为同一对象提供答案时，记录数量不等于独立信息量，
来源置信度也不等于来源正确率。知识获取研究需要把潜在目标、来源机制和查询制度
分别建模，再决定怎样融合证据。

设未知二元目标$Y\in\{0,1\}$，先验概率为
$\mathbb{P}(Y=1)=1/2$。两个来源给出
$L_1,L_2\in\{0,1\}$。先采用一个可检查的简化模型：
给定$Y$后两来源条件独立，且正确率均为$1/2<p<1$，
\[
  \mathbb{P}(L_j=Y\mid Y)=p,
  \qquad j=1,2.
\]
若两来源都报告$1$，Bayes公式给出后验赔率
\begin{align}
  \frac{
    \mathbb{P}(Y=1\mid L_1=L_2=1)
  }{
    \mathbb{P}(Y=0\mid L_1=L_2=1)
  }
  &=
  \frac{\mathbb{P}(Y=1)}{\mathbb{P}(Y=0)}
  \prod_{j=1}^{2}
  \frac{
    \mathbb{P}(L_j=1\mid Y=1)
  }{
    \mathbb{P}(L_j=1\mid Y=0)
  }
  \nonumber\\
  &=
  \left(\frac{p}{1-p}\right)^2.
  \label{eq:future-source-fusion-odds}
\end{align}
这里回用第~\ref{chap:random-variables}章“随机变量”的两层工具：
第~\ref{sec:rv-relations}节在来源规律给定时计算联合、条件与Bayes更新，
第~\ref{sec:stats-targets-data-sources}节则要求在规律未知时先判断
$p$及来源关系能否由观测识别。

\begin{example}[独立来源与复制来源]
  \label{ex:future-source-dependence}
  取$p=0.8$。若条件独立，两次一致报告$1$产生的后验概率为
  \[
    \frac{0.8^2}{0.8^2+0.2^2}
    =
    \frac{16}{17}
    \approx0.941.
  \]
  若第二个来源只是复制第一个来源，那么观察
  $L_1=L_2=1$只相当于一次报告，后验概率仍为$0.8$。
  把复制误写成条件独立，会制造并不存在的证据增益。
\end{example}

在上述严格简化模型中，未知的\(p\)可以由多个对象上的无真值来源输出识别。事实上，
两来源的一致率满足
\[
  a
  =
  \mathbb{P}(L_1=L_2)
  =
  p^2+(1-p)^2.
\]
该函数在\(p>1/2\)上严格递增，因此总体联合分布唯一确定
\[
  p
  =
  \frac{1+\sqrt{2a-1}}{2}.
\]
有限样本只会带来对\(a\)和\(p\)的估计误差，不构成总体层面的不可识别性。

不可识别问题出现在更一般的来源模型中。例如，去掉\(p>1/2\)的方向约束会保留
\(p\)与\(1-p\)的标签翻转；允许两个来源具有不同未知正确率时，一致率不足以分别确定
两个参数；若进一步同时把类别先验和来源正确率设为未知，或允许未知的条件相关结构，
不同参数组合也可能产生相同的来源联合分布。
层次模型可以在多个对象间共享来源质量，但共享先验不会自动创造识别信息。只有在这类
歧义确实存在时，才需要把锚点、已知质量来源、结构约束或少量可靠标签写进模型。
更多来源也只有在其误差结构提供新增信息时才有帮助。

\begin{example}[相同来源输出与不同潜在解释]
  \label{ex:future-source-quality-nonidentifiability}
  固定$P(Y=1)=1/2$，仍假设给定$Y$后来源独立，但允许正确率分别为$p_1,p_2$。
  模型甲取$(p_1,p_2)=(0.8,0.8)$，模型乙取$(0.95,0.7)$。
  两者都给出$P(L_1=L_2)=0.68$，而且完整的来源输出分布也相同：
  \[
    P(1,1)=P(0,0)=0.34,\qquad
    P(1,0)=P(0,1)=0.16.
  \]
  例如$P(1,1)=[p_1p_2+(1-p_1)(1-p_2)]/2$，两组参数代入都为$0.34$。
  但收到两个$1$后的潜在真值概率不同：
  \[
    P_{\text{甲}}(Y=1\mid1,1)=\frac{0.64}{0.68}=\frac{16}{17},
    \qquad
    P_{\text{乙}}(Y=1\mid1,1)=\frac{0.665}{0.68}=\frac{133}{136}.
  \]
  因而即使知道无标签来源输出的整个总体分布，也无法区分这两种质量解释及对应后验。
  这与有限样本估计不准不同：增加同类无标签数据不会消除该歧义；限制两来源正确率
  相同或增加可靠标签，才会改变可识别条件。
\end{example}

\subsection{规则、软标签与关系监督}
\label{sec:future-rules-model-outputs-learning-objectives}

规则来源改变了表示方式。命题逻辑与约束可以排除不允许的答案，
软约束则把违反程度加入优化目标。第~\ref{chap:combinatorics}章提供有限约束、
计数和求解边界，第~\ref{chap:optimization}章说明惩罚、拉格朗日乘子和可行性。
惩罚系数有限时，软约束通常不等于严格满足规则；
规则之间冲突时，也必须规定优先级、松弛方式或不可行处理。

已有模型作为来源时，概率评分与校准评价其输出和实际频率之间的关系。
一个模型报告$0.9$，只表示它给出的概率数值，不表示该来源有$90\%$的真实性。
数据自身提供的增强视图、邻域或对比关系也属于来源假设：
第~\ref{chap:differential-geometry-and-symmetry}章可检查变换关系，
第~\ref{chap:information-theory-and-statistical-geometry}章可比较分布与表示，
但“哪些变化不改变任务语义”仍需由任务定义或证据支持。

\begin{example}[同一对象上的三种监督]
  \label{ex:future-supervision-three-forms}
  设模型对对象\(\bm x_1\)给出的正类概率为\(p_1=0.4\)，对相邻对象
  \(\bm x_2\)给出\(p_2=0.7\)，并以
  \(\widehat y_i=\mathbb{I}\{p_i\geq1/2\}\)形成类别决策。
  人工规则若断言“\(\bm x_1\)必须属于正类”，可以把
  \(\widehat y_1=1\)写成硬约束；当前决策不在可行集合中。若把规则改成有限权重的
  违反代价，模型就可以用其他目标收益抵消违规，不再具有同样的逻辑含义。

  已有模型若为\(\bm x_1\)提供软标签\(q_1=0.8\)，概率监督可以加入交叉熵
  \[
    -q_1\log p_1-(1-q_1)\log(1-p_1)
    \approx0.835.
  \]
  这个目标保留了来源报告的不确定程度，但\(0.8\)是否校准、是否代表来源正确率，
  仍需另行验证。若数据关系只说明\(\bm x_1\)与\(\bm x_2\)应有相同标签，
  一种关系监督是加入
  \[
    (p_1-p_2)^2=0.09.
  \]
  它约束两个输出的一致性，却没有单独说明二者应接近\(0\)还是\(1\)。
  三种信号作用于同一对象，但分别提供可行性、目标分布和对象间关系；
  数值\(0.835\)与\(0.09\)还受损失尺度和权重影响，不能直接当作监督强弱排序。
\end{example}

\subsection{主动查询与信息价值}
\label{sec:future-passive-information-active-query}

主动查询把固定观测改成按历史选择信息。
若当前信息为$\mathcal{F}_t$，查询$q$产生观测$O_q$，
将定义~\ref{def:decision-expected-sample-information-value}用于当前查询，信息价值写为
\[
  \operatorname{VOI}(q)
  =
  \inf_a
  \mathbb{E}[\ell(a,Y)\mid\mathcal{F}_t]
  -
  \mathbb{E}\!\left[
    \inf_a
    \mathbb{E}[\ell(a,Y)\mid\mathcal{F}_t,O_q]
    \,\middle|\,
    \mathcal{F}_t
  \right].
\]
它衡量查询后重新决策能减少多少条件风险。
只要查询后允许忽略$O_q$并沿用原行动，且上述条件期望存在，
第二项的最优风险不会高于第一项，因而$\operatorname{VOI}(q)\geq0$；
若行动集合或损失也随查询改变，就不能直接沿用这个比较。
查询成本为$c(q)$时，只有价值与成本的比较才决定是否查询。
第~\ref{chap:decision-theory-and-dynamic-risk}章提供信息价值，
第~\ref{chap:stochastic-processes-and-approximation}章提供自适应取样后的误差工具。

继续采用正确率$0.8$、条件独立的来源模型，并把行动限制为二元预测、损失取零一损失。
已见$L_1=1$时，当前最优预测为$1$，条件风险为$0.2$。
查询同质量的第二个来源，报告$1$与$0$的条件概率分别为$0.68$与$0.32$。
前一种结果使后验为$16/17$，后一种使后验为$1/2$，所以查询后的平均最优风险为
\[
  0.68\cdot\frac1{17}+0.32\cdot\frac12=0.2.
\]
这次查询虽改变后验，当前决策的信息价值却为零。它不能严格改善两种报告结果下的
平均分类风险；因此任何正查询成本都无法由当前这一项收益抵偿。
若已经收到相互冲突的两个报告，当前后验为$1/2$，再查询一个独立同质量来源，
风险就从$0.5$降到$0.2$，信息价值为$0.3$。
同一个来源是否值得询问，取决于已有信息与当前损失。

若已见$L_1=1$后可改查正确率$0.95$的条件独立来源，它报告$1$的概率为$0.77$，
对应后验为$76/77$；报告$0$的概率为$0.23$，对应后验为$4/23$。
查询后平均风险为$0.77/77+0.23(4/23)=0.05$，价值因而为$0.15$。
成本为$0.1$时值得查询，成本为$0.2$时则不值得。来源更可靠、后验信息增加、
决策风险降低与净收益为正，是四个需要分别判断的命题。

本节推导建立在来源正确率已定义、条件独立结构正确以及先验可用的条件下。
实际知识获取还会遇到选择性缺失、来源共谋、任务漂移与反馈影响。
将信息写入后验、约束或决策之前，必须说明它由什么机制产生、重复记录是否增加证据，
以及最终要改善哪一种损失。来源融合的公式不能代替这些任务判断。

\section{判断旧证据能否用于变化任务}
\label{sec:future-knowledge-evolution-transfer}

前节在一个目标下改变获取证据的方式。本节进一步允许目标分布、任务或环境随使用场景改变。
旧证据能否继续使用，取决于支持、条件规律与表示对应；随时间更新时，还需同时控制
新环境的跟踪误差和旧目标的能力保留。这些条件把概率换测度、几何表示与动态稳定性连接起来。

\subsection{协变量偏移与风险重加权}
\label{sec:future-input-shift-risk-reweighting}

把已有模型用于新环境时，核心问题不是两个参数是否接近，
而是目标环境中的风险能否由已有数据识别和估计。
输入分布变化提供了一个最小例子：若条件标签规律保持，
并且目标分布没有把质量放到来源从未覆盖的区域，
来源样本可以通过密度比重加权目标平均。

设来源联合分布与目标联合分布分别为$P_{\mathrm s}$和$P_{\mathrm t}$，
并假设
\[
  P_{\mathrm t}(Y\mid X)
  =
  P_{\mathrm s}(Y\mid X).
\]
若目标输入分布相对于来源输入分布绝对连续，
密度比
\[
  w(x)
  =
  \frac{\mathrm dP_{\mathrm t}^{X}}
  {\mathrm dP_{\mathrm s}^{X}}(x)
\]
在来源几乎处处有定义。对可积损失，
\begin{align}
  R_{\mathrm t}(h)
  &=
  \mathbb{E}_{P_{\mathrm t}}
  [\ell(h(X),Y)]
  \nonumber\\
  &=
  \mathbb{E}_{P_{\mathrm s}}
  \left[
    w(X)\ell(h(X),Y)
  \right].
  \label{eq:future-covariate-shift}
\end{align}
第~\ref{chap:probability-theory}章的换测度说明等式为何成立，
第~\ref{chap:information-theory-and-statistical-geometry}章的分布差异
则可进一步控制目标变化的大小。

\begin{example}[两点空间上的重加权]
  \label{ex:future-two-point-reweighting}
  令$X\in\{a,b\}$，
  \[
    P_{\mathrm s}^{X}(a)=\frac{3}{4},
    \quad
    P_{\mathrm s}^{X}(b)=\frac{1}{4},
    \qquad
    P_{\mathrm t}^{X}(a)
    =
    P_{\mathrm t}^{X}(b)
    =
    \frac{1}{2}.
  \]
  密度比分别为$w(a)=2/3$和$w(b)=2$。
  若某预测规则在$a,b$处的条件期望损失分别为$0.1,0.5$，
  则来源平均风险为
  \[
    \frac{3}{4}(0.1)+\frac{1}{4}(0.5)=0.2,
  \]
  而重加权来源风险为
  \[
    \frac{3}{4}\frac{2}{3}(0.1)
    +
    \frac{1}{4}2(0.5)
    =
    0.3,
  \]
  恰好等于目标平均风险
  $\frac12(0.1)+\frac12(0.5)$。
\end{example}

重加权等式并不保证估计容易。若$w(X)$很大，
加权损失的方差可能很高；若目标在某区域有正概率而来源概率为零，
绝对连续性失败，式~\eqref{eq:future-covariate-shift}便不能从来源数据恢复
该区域的目标风险。第~\ref{chap:causal-inference}章的部分识别思路
可以在额外边界假设下给出范围，但不能产生点识别。
若$P(Y\mid X)$本身改变，协变量偏移假设也已经失效。

为单独看清权重的波动，在原例中进一步假设损失给定$X$后是确定的：
$\ell(a)=0.1,\ell(b)=0.5$，而不是仅规定这两个条件均值。
加权损失$w(X)\ell(X)$于是以$3/4,1/4$的概率分别取$1/15,1$，
在来源分布下的方差为
\[
  \frac34\left(\frac1{15}\right)^2+\frac14-0.3^2
  =\frac{49}{300}.
\]
$n$个独立来源样本的加权均值方差为$49/(300n)$；
未加权损失均值的方差只有$0.03/n$，却估计的是来源风险$0.2$。
较小方差不能补救目标错误。若实际损失给定$X$仍有随机性，还须加入条件方差项；
若密度比也由数据估计，则还要分析权重估计误差。

表~\ref{tab:future-reweighting-contributions}将两点空间算例按位置拆开。
重加权并未改写每个位置的条件损失，只改变该位置代表目标总体的质量；$b$在来源中
较少见，却在目标中占一半，因而每个来源$b$样本的权重较大。若来源根本没有$b$，
这一步就没有可使用的条件损失信息，权重也不能定义为一个有限数。

\begin{table}[htbp]
  \centering
  \begin{tabular}{lcccc}
    \toprule
    位置 & 来源概率 & 目标概率 & 权重 & 目标风险贡献\\
    \midrule
    $a$ & $3/4$ & $1/2$ & $2/3$ & $0.05$\\
    $b$ & $1/4$ & $1/2$ & $2$ & $0.25$\\
    \bottomrule
  \end{tabular}
  \caption{协变量偏移教学例中各位置的贡献。条件期望损失在$a,b$处分别为$0.1,0.5$，
  来源与目标共享该条件规律。最后一列为来源概率乘权重再乘条件损失，合计为目标风险
  $0.3$；仅靠输入重加权不能处理条件标签规律本身改变的情形。}
  \label{tab:future-reweighting-contributions}
\end{table}

\subsection{表示对应、多任务共享与适应}
\label{sec:future-representation-multitask-adaptation}

参数或表示迁移还需要定义“对应”。两个模型的参数可以因坐标变换、
节点置换或低秩分解的不唯一性而数值不同，却表示同一个函数；
参数接近也可能在病态方向上产生很大预测差异。
第~\ref{chap:matrix-analysis}章的扰动与低秩工具、
第~\ref{chap:differential-geometry-and-symmetry}章的坐标与对称性
帮助建立可比较对象，但任务风险仍是最终要核对的量。

\begin{example}[同一映射的两组坐标表达]
  \label{ex:future-same-map-two-coordinates}
  令\(T:\mathbb{R}^2\to\mathbb{R}^2\)在标准基
  \(\mathcal B=(\bm e_1,\bm e_2)\)下的矩阵为
  \[
    \bm A_{\mathcal B}
    =
    \begin{bmatrix}
      2&0\\
      0&1
    \end{bmatrix}.
  \]
  改用基
  \(\mathcal C=(\bm c_1,\bm c_2)\)，其中
  \(\bm c_1=\bm e_1+\bm e_2\)、\(\bm c_2=\bm e_2\)。
  从\(\mathcal B\)坐标变到\(\mathcal C\)坐标的矩阵及其逆为
  \[
    \bm S_{\mathcal C\leftarrow\mathcal B}
    =
    \begin{bmatrix}
      1&0\\
      -1&1
    \end{bmatrix},
    \qquad
    \bm S_{\mathcal B\leftarrow\mathcal C}
    =
    \begin{bmatrix}
      1&0\\
      1&1
    \end{bmatrix}.
  \]
  第~\ref{chap:linear-algebra}章的换基公式给出
  \[
    \bm A_{\mathcal C}
    =
    \bm S_{\mathcal C\leftarrow\mathcal B}
    \bm A_{\mathcal B}
    \bm S_{\mathcal B\leftarrow\mathcal C}
    =
    \begin{bmatrix}
      2&0\\
      -1&1
    \end{bmatrix}.
  \]
  两个矩阵元素不同，却表示同一个\(T\)。例如
  \(\bm v=\bm c_1=(1,1)^{\top}\)满足
  \[
    [\bm v]_{\mathcal B}
    =
    \begin{bmatrix}1\\1\end{bmatrix},
    \qquad
    [\bm v]_{\mathcal C}
    =
    \begin{bmatrix}1\\0\end{bmatrix}.
  \]
  在标准坐标中\(T(\bm v)=(2,1)^{\top}\)；在\(\mathcal C\)坐标中，
  \[
    \bm A_{\mathcal C}[\bm v]_{\mathcal C}
    =
    \begin{bmatrix}2\\-1\end{bmatrix}
    =
    [T(\bm v)]_{\mathcal C},
  \]
  而\(2\bm c_1-\bm c_2=(2,1)^{\top}\)。因此比较迁移前后的参数数组之前，
  必须先说明坐标或对称代表怎样对应；矩阵数值不相同本身不表示函数关系改变。
\end{example}

多个任务共享信息时，可以用层次模型让任务参数围绕共同结构变化，
也可以用低秩矩阵限制任务间自由度。共享越强，数据少的任务可能获得更低方差，
但真实任务差异若超出共享结构，也会引入偏差。
多目标优化说明任务梯度冲突时不存在同时改善所有目标的无条件更新，
双层优化与隐式求导则组织“用一组数据选择适应规则、用另一组数据评价”的依赖。
这些数学形式仍需要明确任务抽样和数据划分。

以两个无噪声二次任务为例，
$F_1(\theta)=\tfrac12(\theta-1)^2$和$F_2(\theta)=\tfrac12(\theta+1)^2$。
强制两任务使用同一个参数时，平均目标在$\theta=0$达到最小，
每个任务仍损失$1/2$；允许各用自己的参数则可以都达到零。
软共享可改为让任务$j$的参数围绕共同中心$m$，最小化
\[
  \frac12(\theta_j-a_j)^2+\frac{\lambda}{2}(\theta_j-m)^2,
  \qquad a_1=1,\quad a_2=-1,
\]
其解为$\theta_j=(a_j+\lambda m)/(1+\lambda)$。
取$m=0,\lambda=1$，两任务参数为$1/2,-1/2$，各自任务损失降为$1/8$。
这一计算说明共享强度怎样限制适应，不是无噪声任务中共享一定有益的证据。
若用无偏估计值$\widehat a_j$代替$a_j$，其方差为$v_j$且$m$预先固定，
收缩后的参数方差为$v_j/(1+\lambda)^2$，
但偏差为$\lambda(m-a_j)/(1+\lambda)$；任务关系决定这个交换是否值得。

\subsection{变化跟踪与旧能力保留}
\label{sec:future-fixed-task-tracking-retention}

持续学习还增加时间轴。若环境在时刻$t$的最优参数为$\bm\theta_t^\star$，
一个简单的跟踪误差分解可写成
\[
  \lVert
    \bm\theta^{(t+1)}-\bm\theta_{t+1}^\star
  \rVert_2
  \leq
  \lVert
    \bm\theta^{(t+1)}-\bm\theta_t^\star
  \rVert_2
  +
  \lVert
    \bm\theta_t^\star-\bm\theta_{t+1}^\star
  \rVert_2.
\]
第一项反映在当前目标上的更新误差，
第二项反映环境本身的漂移。
固定平均动力学下的收敛结论只控制第一类行为；
若第二项长期不小，算法最多讨论跟踪范围，而不能声称收敛到固定点。

对移动二次目标$F_t(\theta)=\tfrac12(\theta-a_t)^2$，无噪声梯度更新为
$\theta^{(t+1)}=(1-\eta)\theta^{(t)}+\eta a_t$。
若$0<\eta\leq1$且$|a_{t+1}-a_t|\leq d$，记$e_t=\theta^{(t)}-a_t$，便有
\begin{align*}
  |e_{t+1}|&\leq(1-\eta)|e_t|+d,\\
  |e_t|&\leq(1-\eta)^t|e_0|+\frac d\eta\bigl(1-(1-\eta)^t\bigr).
\end{align*}
例如$\eta=0.5,d=0.1$时，长期误差上界为$0.2$；
当目标每步恰好增加$0.1$时，误差确实趋于$-0.2$，不会趋于零。
若再使用带噪梯度，增大步长虽能减轻滞后，却也会改变噪声传播，需另作矩分析。

旧能力保留与新环境适应也不是同一个指标。
保留需要指定哪些旧任务、输出或行为不得退化，
适应则需要指定新任务上的收益；定向修正还要规定哪些内容允许改变。
第~\ref{chap:mathematical-statistics}章的时间抽样、重采样和选择后证据
支持持续评价，但反复按结果调整测试集会削弱证据。

沿用两个二次任务，旧参数$\theta=1$在旧任务上损失为零，
在新任务$F_2$上损失为$2$。对新任务取一步$\eta=1/2$的梯度更新后，
$\theta$变为零，新任务损失降至$1/2$，旧任务损失却升至$1/2$。
“适应成功”与“旧能力不退化”在这个例子中就是两项冲突的要求。
加入旧样本回放、参数约束或多目标权重之前，应先说明允许哪一种退化、如何测量；
约束形式本身不能让冲突消失。

本节的密度比推导只覆盖条件规律保持且支持充分的输入分布变化；
低秩、参数距离或表示对齐只提供结构线索。
能否从旧经验获得目标任务收益，仍取决于任务关系、数据覆盖和评价协议。

\section{本章小结}
\label{sec:mathematics-and-future-research-summary}

章首提出的两个问题现在可以分别作答。训练损失为零、参数更新稳定，只能说明某些
计算或训练目标已经满足；它们既没有把训练数据之外的总体纳入比较，也没有检查控制
轨迹是否越界。前者需要针对数据依赖选择的误差控制，后者需要在状态、行动和扰动约束下
逐步核验安全保持。共同温度例中的误差因子$0.5$与安全区间计算承担的正是不同任务。

本章反复使用相同输入来改变问题。一个矩阵用于回归、秩一压缩和聚类时，
零残差、重构误差与簇内平方和不能互换；一个联合分布用于条件查询、参数学习和生成时，
后验概率、似然与样本分布也不能互换。覆盖率约束随机未来结果，扰动界约束指定集合内的
最坏变化，隐私约束相邻数据集上的输出事件。行动、来源和任务发生变化后，
还要重新说明谁产生数据、什么信息可用、哪些目标可识别以及旧分布是否覆盖新任务。

因此，复用数学结论的可靠顺序是先写出所求对象，再检查比较范围与成立条件。
投影要说明内积和子空间，谱要说明矩阵及其作用方式，固定点要说明映射、
空间和收敛条件。更多样本不能解除不可识别性，更快求解不能补足统计保证，
更强共享也不能消除真实任务冲突。后续模型、范式、应用与系统研究会增加具体的数据
协议和资源约束；这些约束只有进入同一推导，数学结论才真正回答当前研究问题。
